%% Question bank chapter file - generated by bank_tools.py build (do not edit while the bank is still being built:
%% it is regenerated from the per-test data in Question Bank\data). Every question is one bankq block:
%%   \begin{bankq}[id=..., source={...}, paper=.., qno=.., type=..., subtopic={...}, answer={...}, answer_src=key|solved,
%%                 working={...}, difficulty=1-10, diff_parts={K S I C F}, diff_note={...},
%%                 group=<passage id or empty>, also_in={duplicate copies in other tests}]
%% difficulty = min(10, K+S+I+C+F-1): K concepts 1-3, S steps 1-4, I insight 0-2, C computation 0-2, F format 0-1
%%   ...question LaTeX, options, figures (\bfig / \bfigi -> img/)...
%%   \end{bankq}
%% Passages shared by several questions are bankpassage blocks; their questions carry group=<passage id>.
\documentclass[10pt]{article}
\usepackage[a4paper,margin=1.8cm]{geometry}
\usepackage[T1]{fontenc}
\usepackage{amsmath,amssymb,mathtools}
\usepackage{graphicx}
\usepackage[export]{adjustbox}
\usepackage{xcolor,needspace,keyval}
\usepackage{enumitem}
\usepackage{booktabs,array,tabularx}
\usepackage[hidelinks]{hyperref}
\setlength{\parindent}{0pt}
\setlength{\parskip}{3pt}
\makeatletter
\define@key{bq}{id}{\def\bq@id{#1}}\define@key{bq}{source}{\def\bq@src{#1}}\define@key{bq}{paper}{\def\bq@paper{#1}}
\define@key{bq}{qno}{\def\bq@qno{#1}}\define@key{bq}{type}{\def\bq@type{#1}}\define@key{bq}{subtopic}{\def\bq@sub{#1}}
\define@key{bq}{answer}{\def\bq@ans{#1}}\define@key{bq}{answer_src}{\def\bq@asrc{#1}}\define@key{bq}{working}{\def\bq@work{#1}}
\define@key{bq}{group}{\def\bq@group{#1}}\define@key{bq}{also_in}{\def\bq@also{#1}}
\define@key{bq}{difficulty}{\def\bq@diff{#1}}\define@key{bq}{diff_parts}{\def\bq@dparts{#1}}\define@key{bq}{diff_note}{\def\bq@dnote{#1}}
\newenvironment{bankq}[1][]{\def\bq@work{}\def\bq@also{}\def\bq@group{}\def\bq@diff{?}\def\bq@dparts{}\def\bq@dnote{}\setkeys{bq}{#1}%
  \par\needspace{4\baselineskip}\medskip\noindent{\small\textbf{\bq@id}\hfill\color{gray}[\bq@src, Paper \bq@paper, Q\bq@qno]
  \bq@type\ $\cdot$ \bq@sub\ $\cdot$ \textbf{D\bq@diff}}\par\nopagebreak\relax}
  {\par{\small\color{teal}Answer: \textbf{\bq@ans}\ifx\bq@work\@empty\else\ (solved: \bq@work)\fi
  \ifx\bq@also\@empty\else\par\color{gray}Also in: \bq@also\fi}\par}
\define@key{bp}{id}{\def\bp@id{#1}}\define@key{bp}{source}{\def\bp@src{#1}}\define@key{bp}{questions}{\def\bp@qs{#1}}
\newenvironment{bankpassage}[1][]{\setkeys{bp}{#1}\par\needspace{4\baselineskip}\bigskip\noindent
  {\small\textbf{Passage \bp@id}\hfill\color{gray}[\bp@src; questions \bp@qs]}\par\nopagebreak\relax\itshape}{\par}
\makeatother
\newcommand{\bfig}[2]{\par\smallskip{\centering\includegraphics[width=#1pt,max width=0.62\linewidth,max height=0.32\textheight]{img/#2}\par}\smallskip}
\newcommand{\bfigi}[2]{\includegraphics[width=#1pt,max width=\linewidth,max height=0.32\textheight,valign=c]{img/#2}}
\newcommand{\bmissing}[1]{\textit{[figure missing in the source conversion: \detokenize{#1}]}}
\newenvironment{qoptions}{\begin{itemize}[leftmargin=1.8em,itemsep=1pt,topsep=2pt,label={}]}{\end{itemize}}
\newcommand{\opt}[2]{\item[(#1)] #2}
\providecommand{\fwstart}{}\providecommand{\fwend}{}\providecommand{\qfullsplit}{}
\newcolumntype{L}[1]{>{\raggedright\arraybackslash}p{#1}}
\newcolumntype{C}[1]{>{\centering\arraybackslash}p{#1}}
\newcolumntype{R}[1]{>{\raggedleft\arraybackslash}p{#1}}
\providecommand{\sectionhdr}[1]{}\providecommand{\subsectionhdr}[1]{}


\begin{document}
{\Large\bfseries Limits --- Question Bank}\par\smallskip
{\small Mathematics chapter M21. Questions from all converted test papers, grouped by source series. Answers: printed key where it matches the paper, otherwise solved (answer\_src=solved, with working). Exact repeats across papers appear once, with also\_in listing the other copies.}\par\medskip\hrule

\section*{Score-I papers}
\begin{bankq}[id=SI12-P1-MATH-S3-Q08, source={Score-I - Test 12 (09 Mar 2022)}, paper=1, qno=8, type=numerical, subtopic={Standard limits}, answer={3}, answer_src=key, difficulty=4, diff_parts={K2 S1 I1 C0 F1}, diff_note={concepts: limits with [.] and {.}, 1\^{}infinity form | steps: [f] = 1 and {f} = tan x/x - 1 ~ x\^{}2/3; evaluate the limit of the 1\^{}inf form | I: expand tan x}]
Let $f(x)=\dfrac{\tan x}x$ and $\displaystyle\lim_{x\to0}\left([f(x)]+x^2\right)^{\frac1{\{f(x)\}}}=e^\lambda$, then find $\lambda$\\
{(where $[.]$ and $\{.\}$ denotes greatest integer and fractional part function respectively)}
\end{bankq}

\section*{Score-II papers}
\begin{bankq}[id=SII02-P1-MATH-S3-Q06, source={Score-II - Test 02 (28 Mar 2021)}, paper=1, qno=6, type=numerical, subtopic={Standard limits}, answer={3}, answer_src=key, difficulty=8, diff_parts={K3 S1 I2 C2 F1}, diff_note={concepts: standard problem in this chapter, standard limits, apply the relevant results | steps: apply the relevant results; compute}]
The value of limit
\[\lim_{x\to0}\dfrac{\sin(3x+a)-3\sin(2x+a)+3\sin(x+a)-\sin a}{x^3}\]
is $-\cos\left(\dfrac{3a}{\lambda}\right)$. Then value of $\lambda$ is (where $\lambda\in\mathbb{N}$)
\end{bankq}
\begin{bankq}[id=SII02-P2-MATH-S1-Q06, source={Score-II - Test 02 (28 Mar 2021)}, paper=2, qno=6, type=multi, subtopic={Standard limits}, answer={C,D}, answer_src=key, difficulty=8, diff_parts={K3 S1 I2 C2 F1}, diff_note={concepts: standard problem in this chapter, standard limits, apply the relevant results | steps: apply the relevant results; compute}]
The value of
\[\resizebox{0.95\linewidth}{!}{$\displaystyle\lim_{x\to\frac{\pi}{2}}\left(12\tan^2x\left(\sqrt{6+3\sin x-2\cos^2x}-\sqrt{3+6\sin x-\cos^2x}\right)\right)$}\]
is less than:

\begin{qoptions}
\opt{A}{0}
\opt{B}{1}
\opt{C}{2}
\opt{D}{3}
\end{qoptions}
\end{bankq}
\begin{bankq}[id=SII04-P1-MATH-S3-Q01, source={Score-II - Test 04 (07 May 2021)}, paper=1, qno=1, type=numerical, subtopic={Standard limits}, answer={1}, answer_src=key, difficulty=8, diff_parts={K3 S1 I2 C2 F1}, diff_note={concepts: standard problem in this chapter, standard limits, apply the relevant results | steps: apply the relevant results; compute}]
The value of limit $\displaystyle\lim_{x\to0^-}\left[\dfrac{(1+x)^{1/x}}{e}\right]$ is: (Where [.] represent greatest integer function)
\end{bankq}
\begin{bankq}[id=SII08-P1-MATH-S3-Q08, source={Score-II - Test 08 (05 Feb 2023)}, paper=1, qno=8, type=numerical, subtopic={Standard limits}, answer={3.00}, answer_src=key, difficulty=8, diff_parts={K3 S1 I2 C2 F1}, diff_note={concepts: standard problem in this chapter, standard limits, apply the relevant results | steps: apply the relevant results; compute}]
If
$$\beta = \lim_{x\to0}\dfrac{e^{x^3}-(1-x^3)^{1/3}+\left((1-x^2)^{1/2}-1\right)\sin x}{x\sin^2x}$$
then the value of $\dfrac{18}{5}\beta$ is\_\_\_\_\_\_.
\end{bankq}
\begin{bankq}[id=SII08-P2-MATH-S1-Q04, source={Score-II - Test 08 (05 Feb 2023)}, paper=2, qno=4, type=single, subtopic={Standard limits}, answer={B}, answer_src=key, difficulty=4, diff_parts={K2 S1 I1 C1 F0}, diff_note={concepts: standard problem in this chapter, standard limits | steps: apply the relevant results; compute}]
For positive integer n, define\par\noindent
\resizebox{\linewidth}{!}{$f(n) = n + \dfrac{16+5n-3n^2}{4n+3n^2} + \dfrac{32+n-3n^2}{8n+3n^2} + \dfrac{48-3n-3n^2}{12n+3n^2} + \ldots + \dfrac{25n-7n^2}{7n^2}.$}\par
Then the value of $\displaystyle e^{4-\lim_{n\to\infty}f(n)}$ is equal to
\begin{qoptions}
\opt{A}{$\left(\dfrac73\right)^{4/3}$}
\opt{B}{$\left(\dfrac73\right)^{3/4}$}
\opt{C}{$\left(\dfrac74\right)^{4/3}$}
\opt{D}{$\left(\dfrac74\right)^{3/4}$}
\end{qoptions}
\end{bankq}
\begin{bankq}[id=SII09-P2-MATH-S3-Q01, source={Score-II - Test 09 (12 Feb 2023)}, paper=2, qno=1, type=numerical, subtopic={Standard limits}, answer={8}, answer_src=key, difficulty=8, diff_parts={K3 S1 I2 C2 F1}, diff_note={concepts: standard problem in this chapter, standard limits, apply the relevant results | steps: apply the relevant results; compute}]
Let $f(x)=1-|x-2|$ for $1\le x\le3$ and $f(3x)=3f(x)$ for all positive real x. The number of points in [1, 100] where f(x) is non differentiable, is
\end{bankq}
\begin{bankq}[id=SII12-P2-MATH-S2-Q06, source={Score-II - Test 12 (04 Feb 2024)}, paper=2, qno=6, type=multi, subtopic={Standard limits}, answer={A,D}, answer_src=key, difficulty=8, diff_parts={K3 S1 I2 C2 F1}, diff_note={concepts: standard problem in this chapter, standard limits, apply the relevant results | steps: apply the relevant results; compute}]
Let $[.]$ denotes the greatest integer function, $f:(0,1)-\left\{\dfrac13\right\}\to\mathbb R$ be the function defined as $f(x)=[4x]\left(x-\dfrac{1-[x]}{3}\right)^{[2x]}\left(x-\dfrac{[2x]+1}{4}\right)^{3+[x]}$. Then which of the following options is/are Correct?
\begin{qoptions}
\opt{A}{The function f is dis-continuous at exactly two points.}
\opt{B}{There exists exactly two points in $(0,1)$ at which the function f is continuous but \textbf{NOT} differentiable.}
\opt{C}{The function f is \textbf{NOT} differentiable at more than three points in $(0,1)$.}
\opt{D}{The minimum value of f is 0.}
\end{qoptions}
\end{bankq}
\begin{bankq}[id=SII14-P1-MATH-S2-Q10, source={Score-II - Test 14 (06 Feb 2024)}, paper=1, qno=10, type=multi, subtopic={Standard limits}, answer={A,B,C,D}, answer_src=key, difficulty=8, diff_parts={K3 S1 I2 C2 F1}, diff_note={concepts: standard problem in this chapter, standard limits, apply the relevant results | steps: apply the relevant results; compute}]
Let $\displaystyle\lim_{x\to\infty}\sqrt{x^4+Ax^3+Ax^2+x-1}-\sqrt{Cx^4+2x^3+7x^2-4x+1}=L$, then
\begin{qoptions}
\opt{A}{$C\ne1\Rightarrow$ L does not exist}
\opt{B}{$C=1,A\ne2\Rightarrow$ L does not exist}
\opt{C}{$C=1,A=2\Rightarrow L=-\dfrac52$}
\opt{D}{$A=C\Rightarrow$ L does not exist}
\end{qoptions}
\end{bankq}
\begin{bankq}[id=SII21-P2-MATH-S1-Q04, source={Score-II - Test 21 (14 Mar 2024)}, paper=2, qno=4, type=single, subtopic={Standard limits}, answer={A}, answer_src=key, difficulty=4, diff_parts={K2 S1 I1 C1 F0}, diff_note={concepts: standard problem in this chapter, standard limits | steps: apply the relevant results; compute}]
For positive integer n, define
\par\noindent\resizebox{\linewidth}{!}{$\displaystyle f(n) = n + \dfrac{16+5n-5n^2}{4n+5n^2} + \dfrac{32+n-5n^2}{8n+5n^2} + \dfrac{48-3n-5n^2}{12n+5n^2} + \ldots + \dfrac{25n-5n^2}{9n^2}$}\par
Then, the value of $\displaystyle\lim_{n\to\infty} f(n)$ is equal to
\begin{qoptions}
\opt{A}{$4-\dfrac{11}{4}\log_e\left(\dfrac95\right)$}
\opt{B}{$4-\dfrac34\log_e\left(\dfrac75\right)$}
\opt{C}{$4+\dfrac43\log_e\left(\dfrac73\right)$}
\opt{D}{$3+\dfrac{11}{4}\log\dfrac95$}
\end{qoptions}
\end{bankq}
\begin{bankq}[id=SII21-P2-MATH-S3-Q05, source={Score-II - Test 21 (14 Mar 2024)}, paper=2, qno=5, type=numerical, subtopic={Standard limits}, answer={5}, answer_src=key, difficulty=8, diff_parts={K3 S1 I2 C2 F1}, diff_note={concepts: standard problem in this chapter, standard limits, apply the relevant results | steps: apply the relevant results; compute}]
If
\[\beta = \lim_{x\to0} \dfrac{e^{x^3}-(1-x^3)^{1/3}+\left((1-x^2)^{1/2}-1\right)\tan x}{x\tan^2 x}\]
then the value of $6\beta$ is\underline{\hspace{2em}}.
\end{bankq}
\begin{bankq}[id=SII22-P1-MATH-S2-Q08, source={Score-II - Test 22 (14 Mar 2024)}, paper=1, qno=8, type=numerical, subtopic={Standard limits}, answer={34}, answer_src=key, difficulty=8, diff_parts={K3 S1 I2 C2 F1}, diff_note={concepts: standard problem in this chapter, standard limits, apply the relevant results | steps: apply the relevant results; compute}]
Let $\displaystyle\text{Max}_{0\le x\le2}\left\{\dfrac{9-x^2}{5-x}\right\}=\alpha$ and $\displaystyle\text{Min}_{0\le x\le2}\left\{\dfrac{9-x^2}{5-x}\right\}=\beta$\\
If $\displaystyle\int_{\beta-\frac83}^{2\alpha-1}\text{Max}\left\{\dfrac{9-x^2}{5-x},x\right\}dx=\alpha_1+\alpha_2\log_e\left(\dfrac{8}{15}\right)$\\
then $\alpha_1+\alpha_2$ is equal to \_\_\_\_\_
\end{bankq}
\begin{bankq}[id=SII29-P1-MATH-S1-Q02, source={Score-II - Test 29 (26 Mar 2024)}, paper=1, qno=2, type=single, subtopic={Standard limits}, answer={A}, answer_src=key, difficulty=4, diff_parts={K2 S1 I1 C1 F0}, diff_note={concepts: standard problem in this chapter, standard limits | steps: apply the relevant results; compute}]
If $\displaystyle\lim_{x\to k^-}[\sin^{-1}x+\cos^{-1}x+\tan^{-1}x]$ does not exist, then sum of the all possible values of k is.
[Note: [y] denotes greatest integer function less than or equal to y.]

\begin{qoptions}
\opt{A}{$\dfrac{-\sin3}{\sin1\cdot\sin2}$}
\opt{B}{$\dfrac{-\sin2}{\sin1\cdot\sin3}$}
\opt{C}{$\dfrac{-\sin1}{\sin2\cdot\sin3}$}
\opt{D}{$\dfrac{-\sin1\cdot\sin2}{\sin3}$}
\end{qoptions}
\end{bankq}
\begin{bankq}[id=SII42-P1-MATH-S1-Q02, source={Score-II - Test 42 (19 Mar 2025)}, paper=1, qno=2, type=multi, subtopic={Standard limits}, answer={A,B,D}, answer_src=key, difficulty=8, diff_parts={K3 S1 I2 C2 F1}, diff_note={concepts: standard problem in this chapter, standard limits, apply the relevant results | steps: apply the relevant results; compute}]
Let $\left(1+\dfrac1n\right)^{n+x_n}=e$ and $\displaystyle\lim_{n\to\infty}x_n=k$
$$f(x)=\sum_{r=1}^m\left(x^r+x^{-r}\right)^{1/k},\ x\ne\pm1$$
$$\resizebox{\linewidth}{!}{$g(x)=\begin{cases}\displaystyle\lim_{m\to\infty}(f(x)-2m)x^{-2m-2}\cdot(1-x^2), & x\ne1,-1,0\\[4pt] -1 & \text{for }x=\pm1\end{cases}$}$$
then g(x) is:

\begin{qoptions}
\opt{A}{discontinuous at x = $-1$}
\opt{B}{continuous at x = 2}
\opt{C}{limit exist at x = 1}
\opt{D}{limit does not exist at x = 1}
\end{qoptions}
\end{bankq}
\begin{bankq}[id=SII47-P1-MATH-S2-Q01, source={Score-II - Test 47 (30 Mar 2025)}, paper=1, qno=1, type=numerical, subtopic={Standard limits}, answer={8}, answer_src=key, difficulty=8, diff_parts={K3 S1 I2 C2 F1}, diff_note={concepts: standard problem in this chapter, standard limits, apply the relevant results | steps: apply the relevant results; compute}]
The value of $\displaystyle\lim_{x\to2}\dfrac{2^x+2^{3-x}-6}{2^{-\frac{x}{2}}-2^{1-x}}$ is
\end{bankq}

\section*{Score Advanced papers}
\begin{bankq}[id=SA02-P1-MATH-S3-Q04, source={Score Advanced - Test 02 (13 May 2021)}, paper=1, qno=4, type=numerical, subtopic={Standard limits}, answer={0.40}, answer_src=solved, working={$\left(f(3+3/x)/f(3)\right)^x\to e^{3f'(3)/f(3)}=e$, $[L]=2$, $[L]/5=0.4$}, difficulty=4, diff_parts={K1 S2 I1 C0 F1}, diff_note={concepts: limit of the form 1\^{}infinity | steps: exponent x times the change in ln f; equals 3 f'(3)/f(3) = 1; L = e | I: log derivative}, also_in={TOAS - Test 08 (15 May 2022) P1 Q5; TNAS - Test 04 (18 Jun 2023) P2 Q5}]
Let $f : (0, \infty) \to (0, \infty)$ be a function differentiable at 3 and satisfying $f(3) = 3f'(3)$. If the value of $\displaystyle\lim_{x \to \infty}\left(\dfrac{f\left(3 + \frac{3}{x}\right)}{f(3)}\right)^x$ is L then $\dfrac{[L]}{5}$ is (where [.] is greatest integer function)
\end{bankq}
\begin{bankq}[id=SA02-P2-MATH-S3-Q04, source={Score Advanced - Test 02 (13 May 2021)}, paper=2, qno=4, type=single, subtopic={Standard limits}, answer={1}, answer_src=solved, working={S\_n is the mean of the shifts: L=1/2, M=pi/6, L+M=1.02}, difficulty=7, diff_parts={K2 S2 I2 C2 F0}, diff_note={concepts: limit of a product raised to a power, standard limits | steps: reduce to an average; evaluate each limit; sum}]
S$_n$ = $\displaystyle\lim_{x \to \infty}\bigl(\bigl($ $\left(x + \dfrac{1}{\sqrt{4n^2}}\right)$ $\left(x + \dfrac{1}{\sqrt{4n^2 - 1^\alpha}}\right)$ $\left(x + \dfrac{1}{\sqrt{4n^2 - 2^\alpha}}\right)$ $\cdots$ $\left(x + \dfrac{1}{\sqrt{4n^2 - (n-1)^\alpha}}\right)$ $\bigr)^{1/n} - x\bigr)$, n$\in$N\par
and for $\alpha$ = 1, L = $\displaystyle\lim_{n \to \infty}$ (n. S$_n$) and\par
for $\alpha$ = 2, M = $\displaystyle\lim_{n \to \infty}$ (n. S$_n$)\par
then value of [L + M] is equal to\par
[Note : [y] denotes the greatest integer function of y.]
\end{bankq}
\begin{bankq}[id=SA05-P1-MATH-S2-Q06, source={Score Advanced - Test 05 (15 Jul 2021)}, paper=1, qno=6, type=numerical, subtopic={Standard limits}, answer={2.5}, answer_src=key, difficulty=7, diff_parts={K3 S1 I2 C1 F1}, diff_note={concepts: limit to be finite and non-zero, standard limits, c must be 1 and b must be 1 | steps: c must be 1 and b must be 1; evaluate (1-cos t)/t\^{}2}]
If $\displaystyle\lim_{x\to\sqrt{a}}\dfrac{b-\cos(x^2-a)}{(x^2-a)\sin(cx^2-a)} = $ L, is non-zero finite (a$>$0), then value of (L + b + c) is -
\end{bankq}
\begin{bankq}[id=SA05-P1-MATH-S2-Q08, source={Score Advanced - Test 05 (15 Jul 2021)}, paper=1, qno=8, type=numerical, subtopic={L'Hopital}, answer={0.665}, answer_src=key, difficulty=6, diff_parts={K2 S1 I2 C1 F1}, diff_note={concepts: continuity at a point with an integral, l'hopital | steps: numerator ~(2/3)x\^{}3 for small x}]
If $f(x) = \dfrac{\displaystyle\int_0^{x^2}\sin\sqrt{x}\,dx}{x^3}$, x $\geqslant$ 0 is continuous function at x = 0, then value of $f(0)$ is-
\end{bankq}
\begin{bankq}[id=SA07-P1-MATH-S3-Q01, source={Score Advanced - Test 07 (16 Apr 2023)}, paper=1, qno=1, type=numerical, subtopic={Standard limits}, answer={4}, answer_src=key, difficulty=9, diff_parts={K3 S2 I2 C2 F1}, diff_note={concepts: limit of x ln(e(1+1/x)\^{}(1-x)), standard limits, expand ln(1+1/x) | steps: expand ln(1+1/x); simplify; limit}]
Let $\displaystyle\lim_{\text{x}\to\infty}$ x ln$\left(\text{e}\left(1+\dfrac{1}{\text{x}}\right)^{1-\text{x}}\right)$ equals $\dfrac{\text{m}}{\text{n}}$ where m and n are relatively prime positive integer. Find 2m -- n.
\end{bankq}
\begin{bankq}[id=SA08-P2-MATH-S2-Q01, source={Score Advanced - Test 08 (22 Apr 2023)}, paper=2, qno=1, type=numerical, subtopic={Expansions}, answer={0.08}, answer_src=key, difficulty=8, diff_parts={K3 S1 I2 C2 F1}, diff_note={concepts: value of k for continuity from series expansions, expand numerator to x\^{}4, denominator to x\^{}4 | steps: expand numerator to x\^{}4 and denominator to x\^{}4}]
Let $f$(x) = $\begin{cases}\dfrac{\text{e}^{-\text{x}^2/2}-\cos\text{x}}{\text{x}\,\ell\text{n}(1+\text{x})\sin\text{x}(\text{e}^\text{x}-1)} & ;\ \text{x}\neq0\\[1.5ex] \text{k} & ;\ \text{x}=0\end{cases}$ be continuous at x = 0, then k equals -
\end{bankq}
\begin{bankq}[id=SA10-P2-MATH-S2-Q06, source={Score Advanced - Test 10 (06 May 2023)}, paper=2, qno=6, type=multi, subtopic={Standard limits}, answer={A,B}, answer_src=key, difficulty=8, diff_parts={K3 S1 I2 C2 F1}, diff_note={concepts: limit of a ratio of sums as n tends to infinity, standard limits, convert sums to integrals | steps: convert sums to integrals}]
If \par\resizebox{\linewidth}{!}{$\displaystyle\lim_{n\to\infty}\left(\dfrac{1 + \sqrt[3]{2} + \sqrt[3]{3} + \ldots + \sqrt[3]{n}}{n^{7/3}\left(\dfrac{1}{(\alpha n+1)^2} + \dfrac{1}{(\alpha n+2)^2} + \ldots + \dfrac{1}{n^2(\alpha+1)^2}\right)}\right)$}\par $=$ 15 and $\alpha \in \text{R}^+$ then $\alpha$ is equal to :
\begin{qoptions}
\opt{A}{characteristic value of log$_{10}$3$^9$}
\opt{B}{$\displaystyle\lim_{\text{x}\to 0}\left(\dfrac{\sin(4\sin\text{x})}{\text{x}}\right)$}
\opt{C}{number of non-differentiable points of f(x) = $\left|\text{x}^2 - 3\,|\text{x}| + 2\right|$}
\opt{D}{$\left|\text{Re}\left(\sqrt{23 - 10\sqrt{2}\text{i}}\right)\right|$ (where i = $\sqrt{-1}$)}
\end{qoptions}
\end{bankq}
\begin{bankq}[id=SA13-P1-MATH-S2-Q08, source={Score Advanced - Test 13 (13 May 2023)}, paper=1, qno=8, type=numerical, subtopic={Standard limits}, answer={3}, answer_src=key, difficulty=8, diff_parts={K3 S1 I2 C2 F1}, diff_note={concepts: limit of a polynomial quotient, standard limits, derivative of the sum of powers at x=2 | steps: derivative of the sum of powers at x=2}]
If $\displaystyle\lim_{x\to2}\dfrac{2+x+x^2+x^3+\ldots\ldots\ldots\ldots+x^{n-1}-2^n}{x-2}$ = 4097, then the value of [ $\sqrt{n}$ ] equals (where [*] denotes the greatest integer function)\par
\end{bankq}
\begin{bankpassage}[id=SA13-P2-MATH-S3-G11-11, source={Score Advanced - Test 13 (13 May 2023)}, questions={11-11}]
Let f (x) = $\displaystyle\lim_{n\to\infty}n\left(\text{x}^{\frac{2022}{n}}-\text{x}^{\frac{2023}{n}}\right)$ ; x $>$ 0.\par
\end{bankpassage}
\begin{bankq}[id=SA13-P2-MATH-S3-Q11, source={Score Advanced - Test 13 (13 May 2023)}, paper=2, qno=11, type=single, subtopic={Standard limits}, answer={A}, answer_src=key, difficulty=7, diff_parts={K3 S1 I2 C2 F0}, diff_note={concepts: matching limits and values for a limit function, standard limits, f(x)=ln x | steps: f(x)=ln x}, group=SA13-P2-MATH-S3-G11-11]
Match List-I with List-II and select the correct answer using the code given below the list.\par
{\footnotesize\setlength{\tabcolsep}{1.5pt}\noindent\begin{tabular}{|c|L{3.6cm}|c|c|}
\hline
 & List-I & & List-II \\
\hline
(a) & Number of solution of equation f(|x|) = e$^{-2023x^2}$ & (P) & 1 \\
\hline
(b) & Number of solution of equation f(|x|) = e$^{2023x^2}$ & (Q) & 2 \\
\hline
(c) & Number of solution of equation f(x) = e$^{2022}$ & (R) & 3 \\
\hline
(d) & Number of solution of equation f(x) = e$^{-2022}$ & (S) & 0 \\
\hline
\end{tabular}}\par
\begin{qoptions}
\opt{A}{a $\rightarrow$ Q; b $\rightarrow$ Q; c $\rightarrow$ P; d $\rightarrow$ Q}
\opt{B}{a $\rightarrow$ P; b $\rightarrow$ P; c $\rightarrow$ Q; d $\rightarrow$ S}
\opt{C}{a $\rightarrow$ P; b $\rightarrow$ Q; c $\rightarrow$ P; d $\rightarrow$ Q}
\opt{D}{a $\rightarrow$ Q; b $\rightarrow$ P; c $\rightarrow$ Q; d $\rightarrow$ S}
\end{qoptions}
\end{bankq}
\begin{bankq}[id=SA15-P1-MATH-S1-Q03, source={Score Advanced - Test 15 (18 May 2023)}, paper=1, qno=3, type=single, subtopic={Expansions}, answer={D}, answer_src=solved, working={The limit is -2e\^{}2}, difficulty=7, diff_parts={K3 S1 I2 C2 F0}, diff_note={concepts: limit with exponential and logarithmic expansions, numerator ~ -e\^{}2 x, denominator ~ x/2 | steps: numerator ~ -e\^{}2 x; denominator ~ x/2}, also_in={Score Advanced - Test 32 (08 May 2024) P1 Q3}]
$\displaystyle\lim_{x\to0^+}\dfrac{(1+\tan x)^{2\cot x}-e^{2\cos x}}{1-\cos\left(\ln\left(1+\sin\sqrt x\right)\right)}$
\begin{qoptions}
\opt{A}{$2e^2$}
\opt{B}{$2e$}
\opt{C}{$-2e$}
\opt{D}{$-2e^2$}
\end{qoptions}
\end{bankq}
\begin{bankq}[id=SA17-P2-MATH-S1-Q01, source={Score Advanced - Test 17 (25 May 2023)}, paper=2, qno=1, type=single, subtopic={Standard limits}, answer={A}, answer_src=key, difficulty=7, diff_parts={K3 S1 I2 C2 F0}, diff_note={concepts: limit of log of a product of secants, standard limits, expand ln sec(x) for small x | steps: expand ln sec(x) for small x; sum of squares}, also_in={Score Advanced - Test 35 (15 May 2024) P2 Q1}]
$\displaystyle\lim_{x\to0}\frac{\ln\left(\sec(e^x)\sec(e^2x)\ldots\sec(e^{50}x)\right)}{e^2-e^{2\cos x}}$ is equal to
\begin{qoptions}
\opt{A}{$\frac{e^{100}-1}{2(e^2-1)}$}
\opt{B}{$\frac{e^{100}-1}{(e^2-1)}$}
\opt{C}{$\frac{e^{100}-1}{2e(e^2-1)}$}
\opt{D}{$\frac{e^{50}-1}{2e(e^2-1)}$}
\end{qoptions}
\end{bankq}
\begin{bankq}[id=SA18-P2-MATH-S2-Q07, source={Score Advanced - Test 18 (27 May 2023)}, paper=2, qno=7, type=multi, subtopic={Standard limits}, answer={A,C,D}, answer_src=key, difficulty=8, diff_parts={K3 S1 I2 C2 F1}, diff_note={concepts: limit of an integral expression, standard limits, substitute and expand to leading order | steps: substitute and expand to leading order}]
$\displaystyle\lim_{x\to0}\left(\frac{1}{x^{2022}}\int_{x^{2023}}^{x^{2022}}\left(\frac{(1+t)^{\frac1{2021}}-1}{t}\right)dt\right)=\frac1P$\par
Then
\begin{qoptions}
\opt{A}{Number of positive divisors of $P$ is 4}
\opt{B}{$\ell n\,x=Px$ has exactly two real solutions}
\opt{C}{$2018^x+2019^x+2020^x=P^x$ has exactly one real root}
\opt{D}{Unit digit of $P^P$ is 1}
\end{qoptions}
\end{bankq}
\begin{bankq}[id=SA19-P1-MATH-S1-Q03, source={Score Advanced - Test 19 (04 Apr 2024)}, paper=1, qno=3, type=single, subtopic={Standard limits}, answer={B}, answer_src=key, difficulty=7, diff_parts={K3 S1 I2 C2 F0}, diff_note={concepts: limit of a scaled integral, standard limits, substitute u=ax | steps: substitute u=ax; expansion near 0}]
$\displaystyle\lim_{a\to\infty}a^2\int_0^{1/a^2}\dfrac{\cos^{-1}ax}{\sec^{-1}(ax+2)}dx$ is equal to
\begin{qoptions}
\opt{A}{3}
\opt{B}{$\dfrac32$}
\opt{C}{$\dfrac23$}
\opt{D}{2}
\end{qoptions}
\end{bankq}
\begin{bankq}[id=SA19-P1-MATH-S3-Q04, source={Score Advanced - Test 19 (04 Apr 2024)}, paper=1, qno=4, type=numerical, subtopic={Expansions}, answer={5}, answer_src=key, difficulty=8, diff_parts={K3 S1 I2 C2 F1}, diff_note={concepts: limit using logarithmic expansion, expansions, take the logarithm | steps: take the logarithm; expand to fourth order}]
If the value of $\displaystyle\lim_{x\to\infty}x^4\left(\left(\dfrac{x^2+1}{x^2-1}\right)^{x^2}-e^2\right)$ is equal to $\dfrac{ae^2}{b}$ (where $a$ and $b$ are co - prime) then the value of $\left|a^2-b^2\right|$ is
\end{bankq}
\begin{bankpassage}[id=SA21-P1-MATH-S3-G1-2, source={Score Advanced - Test 21 (08 Apr 2024)}, questions={1-2}]
$A_n$ and $B_n$ be square matrices of order 3 which are defined as $A_n=[a_{ij}]$ and $B_n=[b_{ij}]$ where $a_{ij}=\dfrac{2i+j}{3^{2n}}$ and $b_{ij}=\dfrac{3i-j}{2^{2n}}$ $\forall\,i,j,1\le i,\ j\le3$. If $L=\displaystyle\lim_{n\to\infty}\mathrm{Tr}\left(3A_1+3^2A_2+3^3.A_3+\ldots..+3^n.A_n\right)$,\par
$M=\displaystyle\lim_{n\to\infty}\mathrm{Tr}\left(2B_1+2^2B_2+2^3.B_3+\ldots..+2^n.B_n\right)$\par
(Where Tr (P) = Trace of matrix P). Let $p,n\in N$,\par
$f(x)=1-x^p$ and\par $g_n(x)=\dfrac{n}{\dfrac{1}{f(x)}+\dfrac{1}{f(2x)}+\ldots\ldots+\dfrac{1}{f(nx)}}$ and\par $R=\displaystyle\lim_{x\to0}\dfrac{1-g_n(x)}{x^p}$\par
at $n=5$ and $p=3$
\end{bankpassage}
\begin{bankq}[id=SA21-P1-MATH-S3-Q01, source={Score Advanced - Test 21 (08 Apr 2024)}, paper=1, qno=1, type=numerical, subtopic={Standard limits}, answer={55}, answer_src=key, difficulty=8, diff_parts={K3 S1 I2 C2 F1}, diff_note={concepts: limit of a harmonic type expression, standard limits, expand 1/(1-k x\^{}p) | steps: expand 1/(1-k x\^{}p); limit}, group=SA21-P1-MATH-S3-G1-2]
The value of $100-R$ is equal to \underline{\hspace{1.5cm}}
\end{bankq}
\begin{bankq}[id=SA21-P1-MATH-S3-Q02, source={Score Advanced - Test 21 (08 Apr 2024)}, paper=1, qno=2, type=numerical, subtopic={Standard limits}, answer={66}, answer_src=key, difficulty=8, diff_parts={K3 S1 I2 C2 F1}, diff_note={concepts: sum of limits of traces and a limit, standard limits, geometric series of traces | steps: geometric series of traces; R}, group=SA21-P1-MATH-S3-G1-2]
The value of $L+M+R$ is equal to
\end{bankq}
\begin{bankq}[id=SA25-P2-MATH-S2-Q06, source={Score Advanced - Test 25 (21 Apr 2024)}, paper=2, qno=6, type=multi, subtopic={Standard limits}, answer={B,D}, answer_src=key, difficulty=8, diff_parts={K3 S1 I2 C2 F1}, diff_note={concepts: limit of an infinite product of fractions, standard limits, telescoping product | steps: telescoping product}, also_in={Score Advanced - Test 26 (21 Apr 2024) P2 Q6}]
Consider f : $\left[0,\dfrac\pi2\right]-\left\{\dfrac\pi4\right\}\rightarrow$ R. Let $f(x)=\displaystyle\lim_{n\to\infty}\prod_{k=0}^n\dfrac{\left((\cot x)^{2^k}+1\right)^2}{\left((\cot x)^{2^{k+1}}+1\right)}$, $\forall x\in\left(0,\dfrac\pi2\right]-\left\{\dfrac\pi4\right\}$ and f$\left(\dfrac\pi4\right)$ = 2 and f(0) = 0, Then which of the following is true?
\begin{qoptions}
\opt{A}{f(x) is increasing in $\left(0,\dfrac\pi4\right)$ \& $\left(\dfrac\pi4,\dfrac\pi2\right)$}
\opt{B}{f(x) is decreasing in $\left(\dfrac\pi4,\dfrac\pi2\right)$}
\opt{C}{f(x) has one local maxima}
\opt{D}{f(x) has three local minima}
\end{qoptions}
\end{bankq}
\begin{bankq}[id=SA27-P1-MATH-S2-Q07, source={Score Advanced - Test 27 (24 Apr 2024)}, paper=1, qno=7, type=multi, subtopic={Standard limits}, answer={A,B}, answer_src=key, difficulty=8, diff_parts={K3 S1 I2 C2 F1}, diff_note={concepts: limit of a function with a given ratio, standard limits, series expansion | steps: series expansion}]
f(x) is the function such that $\displaystyle\lim_{x\to0}\dfrac{\mathrm f(\mathrm x)}{\mathrm x}=1$ if $\displaystyle\lim_{x\to0}\dfrac{\mathrm x(1+\mathrm a\cos\mathrm x)-\mathrm b\sin\mathrm x}{(\mathrm f(\mathrm x))^3}=1$ then
\begin{qoptions}
\opt{A}{a = $\dfrac{-5}{2}$}
\opt{B}{b = $\dfrac{-3}{2}$}
\opt{C}{a = $\dfrac{-7}{2}$}
\opt{D}{b = $\dfrac{-5}{2}$}
\end{qoptions}
\end{bankq}
\begin{bankq}[id=SA28-P1-MATH-S3-Q06, source={Score Advanced - Test 28 (27 Apr 2024)}, paper=1, qno=6, type=numerical, subtopic={Standard limits}, answer={1}, answer_src=key, difficulty=8, diff_parts={K3 S1 I2 C2 F1}, diff_note={concepts: limit involving means of exponentials, standard limits, logarithm and expansions | steps: logarithm and expansions}]
$\displaystyle\lim_{x\to0}\dfrac{\left(\frac{\mathrm a^{\mathrm x}+\mathrm b^{\mathrm x}+\mathrm c^{\mathrm x}}{3}\right)^{2/\mathrm x}}{(\mathrm{abc})^{2/3}}$ is :-\par
(Where a, b, c $>$ 0)
\end{bankq}
\begin{bankq}[id=SA29-P1-MATH-S3-Q08, source={Score Advanced - Test 29 (01 May 2024)}, paper=1, qno=8, type=numerical, subtopic={Standard limits}, answer={4}, answer_src=key, difficulty=8, diff_parts={K3 S1 I2 C2 F1}, diff_note={concepts: duplicate or standard problem in this chapter, standard limits, apply the relevant results | steps: apply the relevant results; compute}]
If $\displaystyle\lim_{n\to\infty}\left(\dfrac{{}^{3n}\mathrm c_n}{{}^{2n}\mathrm c_n}\right)^{\frac1n}=\dfrac ab$ where a, b are relatively prime positive integers then $\dfrac{a+b+1}{a-b}$ =
\end{bankq}
\begin{bankq}[id=SA30-P2-MATH-S1-Q07, source={Score Advanced - Test 30 (04 May 2024)}, paper=2, qno=7, type=multi, subtopic={Expansions}, answer={A,B,C}, answer_src=key, difficulty=8, diff_parts={K3 S1 I2 C2 F1}, diff_note={concepts: standard problem in this chapter, expansions, apply the relevant results | steps: apply the relevant results; compute}, also_in={Score Advanced - Test 31 (04 May 2024) P2 Q7}]
Consider the limit $\displaystyle\lim_{x\to0}\dfrac{a\tan x+b\tan^{-1}x+c\sin^{-1}x}{x^5}$ = 3, where a, b, c $\in$ R, then correct option/s is/are
\begin{qoptions}
\opt{A}{a = 30}
\opt{B}{b = 10}
\opt{C}{c = $-$40}
\opt{D}{b = $-$5}
\end{qoptions}
\end{bankq}
\begin{bankq}[id=SA32-P2-MATH-S4-Q08, source={Score Advanced - Test 32 (08 May 2024)}, paper=2, qno=8, type=single, subtopic={Standard limits}, answer={8}, answer_src=key, difficulty=7, diff_parts={K3 S1 I2 C2 F0}, diff_note={concepts: standard problem in this chapter, standard limits, apply the relevant results | steps: apply the relevant results; compute}]
Let f(x) = $\displaystyle\lim_{x\to\infty}\dfrac{x^{2n-1}+ax^2+bx}{x^{2n}+1}$. If f(x) is continuous for all x $\in$ R, then the value of a + 8b is
\end{bankq}
\begin{bankq}[id=SA38-P2-MATH-S1-Q03, source={Score Advanced - Test 38 (13 Apr 2025)}, paper=2, qno=3, type=single, subtopic={Standard limits}, answer={B}, answer_src=key, difficulty=4, diff_parts={K2 S1 I1 C1 F0}, diff_note={concepts: standard problem in this chapter, standard limits | steps: apply the relevant results; compute}]
Let k $\in\mathbb R^+$. If $\displaystyle\lim_{x\to0^+}\left(\sin\left(\sin\mathrm k^2\mathrm x\right)+\cos\mathrm x-\mathrm x\right)^{\frac2{\mathrm x}}=\mathrm e^6$, then the value of k is
\begin{qoptions}
\opt{A}{1}
\opt{B}{2}
\opt{C}{3}
\opt{D}{4}
\end{qoptions}
\end{bankq}
\begin{bankq}[id=SA38-P2-MATH-S2-Q05, source={Score Advanced - Test 38 (13 Apr 2025)}, paper=2, qno=5, type=multi, subtopic={Standard limits}, answer={A,B,C,D}, answer_src=key, difficulty=8, diff_parts={K3 S1 I2 C2 F1}, diff_note={concepts: standard problem in this chapter, standard limits, apply the relevant results | steps: apply the relevant results; compute}]
Let S be the set of all ($\alpha$, $\beta$) $\in\mathbb R\times\mathbb R$ such that $\displaystyle\lim_{x\to\infty}\dfrac{\sin\left(x^2\right)\left(\log_ex\right)^\alpha\tan\left(\frac1{x^3}\right)}{x^{\alpha\beta}\left(\log_e(1+x)\right)^\beta}=0$. Then which of the following is (are) correct ?
\begin{qoptions}
\opt{A}{($-$1, 3) $\in$ S}
\opt{B}{($-$1, 1) $\in$ S}
\opt{C}{(1, $-$1) $\in$ S}
\opt{D}{(1, $-$2) $\in$ S}
\end{qoptions}
\end{bankq}
\begin{bankq}[id=SA39-P1-MATH-S1-Q03, source={Score Advanced - Test 39 (20 Apr 2025)}, paper=1, qno=3, type=single, subtopic={Standard limits}, answer={C}, answer_src=key, difficulty=4, diff_parts={K2 S1 I1 C1 F0}, diff_note={concepts: standard problem in this chapter, standard limits | steps: apply the relevant results; compute}]
$\displaystyle\lim_{x\to0}\dfrac{x\mathrm e^{\sin x}-\mathrm e^x\sin^{-1}(\sin x)}{\sin^2x-x\sin x}$ =
\begin{qoptions}
\opt{A}{$-2$}
\opt{B}{$-1$}
\opt{C}{1}
\opt{D}{2}
\end{qoptions}
\end{bankq}
\begin{bankq}[id=SA40-P1-MATH-S4-Q03, source={Score Advanced - Test 40 (27 Apr 2025)}, paper=1, qno=3, type=numerical, subtopic={Standard limits}, answer={3}, answer_src=key, difficulty=8, diff_parts={K3 S1 I2 C2 F1}, diff_note={concepts: standard problem in this chapter, standard limits, apply the relevant results | steps: apply the relevant results; compute}]
Number of solution of the equation log$_\mathrm k$x = cos($\pi$x) is 3 , then sum of all possible integers value of `k' is :
\end{bankq}
\begin{bankq}[id=SA41-P2-MATH-S2-Q06, source={Score Advanced - Test 41}, paper=2, qno=6, type=multi, subtopic={Standard limits}, answer={C,D}, answer_src=key, difficulty=8, diff_parts={K3 S1 I2 C2 F1}, diff_note={concepts: standard problem in this chapter, standard limits, apply the relevant results | steps: apply the relevant results; compute}]
$\displaystyle\lim_{\mathrm n\to\infty}\sum_{\mathrm i=1}^{\mathrm n}\dfrac{\mathrm i}{\mathrm n^2+\mathrm i^2}$ is equal to -
\begin{qoptions}
\opt{A}{$\dfrac14-\ell\mathrm n2$}
\opt{B}{$\dfrac12-\ell\mathrm n2$}
\opt{C}{$\ell\mathrm n\sqrt2$}
\opt{D}{$\dfrac12\ell\mathrm n2$}
\end{qoptions}
\end{bankq}

\section*{TOAS papers}
\begin{bankq}[id=TOAS04-P1-MATH-S1-Q01, source={TOAS - Test 04 (15 Mar 2022)}, paper=1, qno=1, type=single, subtopic={Standard limits}, answer={B}, answer_src=key, difficulty=4, diff_parts={K2 S2 I0 C1 F0}, diff_note={concepts: standard limits, expansions | steps: log(1 + x) - log 2 near x = 1; the exponential and sine factors; the denominator radicals; combine}]
Value of $\displaystyle\lim_{x\to1}\dfrac{(\log(1+x)-\log2)(3\cdot4^{x-1}-3x)}{\left((7+x)^{1/3}-(1+3x)^{1/2}\right)\sin\pi x}$ is equal to
\begin{qoptions}
\opt{A}{$\dfrac9\pi(2\log2-1)$}
\opt{B}{$\dfrac9{4\pi}[\log4-1]$}
\opt{C}{$\dfrac9{2\pi}\left[\log4-\dfrac12\right]$}
\opt{D}{$\dfrac2{3\pi}[2\log2-1]$}
\end{qoptions}
\end{bankq}
\begin{bankq}[id=TOAS04-P1-MATH-S1-Q02, source={TOAS - Test 04 (15 Mar 2022)}, paper=1, qno=2, type=single, subtopic={Standard limits}, answer={B}, answer_src=key, difficulty=5, diff_parts={K2 S2 I1 C1 F0}, diff_note={concepts: complex numbers, sums of trigonometric functions | steps: sum of the unit vectors equals 2 times a unit vector; pairwise cosine sum lambda; evaluate the limit with radicals | I: complex numbers}]
If $\theta_1,\theta_2,\theta_3$ are three distinct real values such that \[\dfrac{\sin\theta_1+\sin\theta_2+\sin\theta_3}{\sin(\theta_1+\theta_2+\theta_3)}=\dfrac{\cos\theta_1+\cos\theta_2+\cos\theta_3}{\cos(\theta_1+\theta_2+\theta_3)}=2\] and $\cos(\theta_1+\theta_2)+\cos(\theta_2+\theta_3)+\cos(\theta_3+\theta_1)=\lambda$, then value of \[\lim_{x\to\lambda}\dfrac{\sqrt{x^2-\lambda^2}}{\sqrt{x-\lambda}+\sqrt{x-\sqrt\lambda}}\] is
\begin{qoptions}
\opt{A}{1}
\opt{B}{2}
\opt{C}{3}
\opt{D}{5}
\end{qoptions}
\end{bankq}
\begin{bankq}[id=TOAS04-P1-MATH-S1-Q03, source={TOAS - Test 04 (15 Mar 2022)}, paper=1, qno=3, type=single, subtopic={Standard limits}, answer={C}, answer_src=key, difficulty=4, diff_parts={K1 S2 I1 C1 F0}, diff_note={concepts: functional equation by substitution | steps: replace x by pi/2 - x; solve for f(sin x); take the limit as sin x tends to 1 | I: swap sin and cos}]
If $2f(\sin x)+\sqrt2f(\cos x)=\tan x$, then $\displaystyle\lim_{x\to1^-}(\sqrt{1-x})f((x))$ is - (Where $x\in\mathbb R-\left\{(2n+1)\dfrac\pi2;\ n\in I\right\}$)
\begin{qoptions}
\opt{A}{$\dfrac1{2\sqrt2}$}
\opt{B}{$\sqrt2$}
\opt{C}{$\dfrac1{\sqrt2}$}
\opt{D}{$\dfrac{\sqrt2}3$}
\end{qoptions}
\end{bankq}
\begin{bankq}[id=TOAS04-P1-MATH-S2-Q05, source={TOAS - Test 04 (15 Mar 2022)}, paper=1, qno=5, type=multi, subtopic={Standard limits}, answer={A,B}, answer_src=key, difficulty=5, diff_parts={K2 S2 I1 C0 F1}, diff_note={concepts: limits of the form 1\^{}infinity, fixed points | steps: l = e\^{}(-1) from tan expansion; limit of the rational power is 0; limit of f\_n is 3 | I: fixed point of f\_1}]
If $\displaystyle\lim_{x\to0}\left(\dfrac{\tan x}x\right)^{\frac{\tan x}{x-\tan x}}$ exists and equal to $\ell$, then which of the following is (are) correct ?
\begin{qoptions}
\opt{A}{$\displaystyle\lim_{x\to\infty}\left(\dfrac{2x^2+1}{4x^2-1}\right)^{\frac{x^2}{2+x}}$ is less than $\ell$}
\opt{B}{If $f_1(x)=\dfrac x3+2\ \forall\ x\in\mathbb R$ and $f_n(x)$ is defined by $f_n(x)=f_1\left(f_{n-1}(x)\right)\ \forall\ n\geqslant2$ and $n\in\mathbb N$ then $\displaystyle\lim_{n\to\infty}f_n(x)$ is more than $\ell$}
\opt{C}{$\displaystyle\lim_{x\to\infty}\left(\dfrac{2x^2+1}{4x^2-1}\right)^{\frac{x^2}{2+x}}$ is more than $\ell$}
\opt{D}{If $f_1(x)=\dfrac x3+2\ \forall\ x\in\mathbb R$ and $f_n(x)$ is defined by $f_n(x)=f_1(f_{n-1}(x))\ \forall\ n\geq2$ and $n\in\mathbb N$ then $\displaystyle\lim_{n\to\infty}f_n(x)$ is less than $\ell$}
\end{qoptions}
\end{bankq}
\begin{bankq}[id=TOAS04-P1-MATH-S2-Q08, source={TOAS - Test 04 (15 Mar 2022)}, paper=1, qno=8, type=multi, subtopic={Standard limits}, answer={A,B,C,D}, answer_src=key, difficulty=4, diff_parts={K2 S2 I0 C0 F1}, diff_note={concepts: right triangle trigonometry, limits | steps: express y and x with theta; identity for y\^{}n - x\^{}n; limits of y - x, y\^{}2 - x\^{}2, y\^{}3 - x\^{}3}, also_in={TEAS - Test 01 (16 Apr 2023) P1 Q4}]
A right angled triangle has one leg of length 1, another of length `$x$' \& hypotenuse of length `$y$'. The angle oppsoite to the side $x$ is $\theta$ then -
\begin{qoptions}
\opt{A}{$y^n-x^n=\dfrac{1-\sin^n\theta}{\cos^n\theta}$}
\opt{B}{$\displaystyle\lim_{\theta\to\frac\pi2}(y-x)=0$}
\opt{C}{$\displaystyle\lim_{\theta\to\frac\pi2}\left(y^2-x^2\right)=1$}
\opt{D}{$\displaystyle\lim_{\theta\to\frac\pi2}\left(y^3-x^3\right)$ does not exist.}
\end{qoptions}
\end{bankq}
\begin{bankpassage}[id=TOAS04-P1-MATH-S3-G5-6, source={TOAS - Test 04 (15 Mar 2022)}, questions={5-6}]
Let $a_1=1$ and $a_n=\sin(a_{n-1})$, $n>1$ \& $n\in\mathbb N$. If \[\lim_{n\to\infty}\dfrac{2^{2a_n}-2^{1+a_n}\cdot3^{a_n}+3^{2a_n}}{\cos a_n+1-e^{a_n}-e^{-a_n}}=-\dfrac\alpha3\,\ell n^{\beta}\left(\dfrac23\right),\] (where $\alpha$ and $\beta$ are prime numbers).
\end{bankpassage}
\begin{bankq}[id=TOAS04-P1-MATH-S3-Q05, source={TOAS - Test 04 (15 Mar 2022)}, paper=1, qno=5, type=numerical, subtopic={Expansions}, answer={6.00}, answer_src=key, difficulty=3, diff_parts={K1 S1 I1 C0 F1}, diff_note={concepts: expansion coefficients | steps: expand the given function; match the coefficient | I: series matching}, group=TOAS04-P1-MATH-S3-G5-6]
The value of $3\alpha$ is
\end{bankq}
\begin{bankq}[id=TOAS04-P1-MATH-S3-Q06, source={TOAS - Test 04 (15 Mar 2022)}, paper=1, qno=6, type=numerical, subtopic={Expansions}, answer={4.00}, answer_src=key, difficulty=2, diff_parts={K1 S1 I0 C0 F1}, diff_note={concepts: expansion coefficients | steps: expand the given function; match the coefficient}, group=TOAS04-P1-MATH-S3-G5-6]
The value of $2\beta$ is
\end{bankq}
\begin{bankq}[id=TOAS04-P1-MATH-S4-Q07, source={TOAS - Test 04 (15 Mar 2022)}, paper=1, qno=7, type=numerical, subtopic={Standard limits}, answer={1}, answer_src=key, difficulty=6, diff_parts={K2 S2 I2 C0 F1}, diff_note={concepts: recurrence and telescoping product, e as a limit | steps: 1 + 1/a\_k = a\_(k+1)/((k+1) a\_k); product telescopes; a\_n/n! tends to e; take the logarithm | I: telescoping product}]
If a sequence is such that $a_1=1$ and $a_n=n(1+a_{n-1})\ \forall\ n\geq2$ and $n\in\mathbb N$, then \[\lim_{n\to\infty}\ell n\left(\left(1+\dfrac1{a_1}\right)\left(1+\dfrac1{a_2}\right)\ldots\left(1+\dfrac1{a_n}\right)\right)\] is \underline{\hspace{2.5em}}.
\end{bankq}
\begin{bankq}[id=TOAS04-P1-MATH-S4-Q09, source={TOAS - Test 04 (15 Mar 2022)}, paper=1, qno=9, type=numerical, subtopic={Expansions}, answer={3}, answer_src=key, difficulty=3, diff_parts={K1 S2 I0 C0 F1}, diff_note={concepts: finite limit forces coefficients | steps: expand log(1 + x) and sin x; a = 0; b = -1; l = -1/2}]
If $\displaystyle\lim_{x\to0}\dfrac{a+\log_e(1+x)+b\sin x}{x^2}=\ell$ (finite) then $|a+b+4\ell|$ (where $a,b,\ell\in\mathbb R$)
\end{bankq}
\begin{bankq}[id=TOAS04-P2-MATH-S1-Q01, source={TOAS - Test 04 (15 Mar 2022)}, paper=2, qno=1, type=multi, subtopic={Standard limits}, answer={C}, answer_src=key, difficulty=5, diff_parts={K2 S2 I1 C0 F1}, diff_note={concepts: ratio of conjugate terms, standard limits | steps: A\_n/B\_n tends to sqrt2; evaluate each option; identify the matching limit | I: ratio of conjugate powers}]
Let $(1+\sqrt2)^n=A_n+B_n\sqrt2$, where $A_n$ and $B_n$ are rational numbers and $n\in N$, then the value of $\displaystyle\lim_{n\to\infty}\dfrac{A_n}{B_n}$ is equal to -\\
(where $[.]$ denotes greatest integer function)
\begin{qoptions}
\opt{A}{$\left[\displaystyle\lim_{x\to0}\dfrac{\sin x}{x}\right]$}
\opt{B}{$\displaystyle\lim_{x\to0}\left[\dfrac{\sin x}{x}\right]$}
\opt{C}{$\displaystyle\lim_{x\to\pi/4}\dfrac{|\cos x-\sin x|}{\left|x-\frac\pi4\right|}$}
\opt{D}{$\left[\displaystyle\lim_{x\to0}\dfrac{\sin\sqrt2\,x}{x}\right]$}
\end{qoptions}
\end{bankq}
\begin{bankq}[id=TOAS04-P2-MATH-S1-Q02, source={TOAS - Test 04 (15 Mar 2022)}, paper=2, qno=2, type=multi, subtopic={Standard limits}, answer={C}, answer_src=key, difficulty=3, diff_parts={K1 S1 I1 C0 F1}, diff_note={concepts: limit at minus infinity with a radical | steps: factor  | x}]
$\displaystyle\lim_{x\to-\infty}\left(\sqrt{x^2-\dfrac{4x}{3}}+x+2\right)$ is equal to-
\begin{qoptions}
\opt{A}{$\dfrac43$}
\opt{B}{$-\dfrac43$}
\opt{C}{$\dfrac83$}
\opt{D}{$-\dfrac83$}
\end{qoptions}
\end{bankq}
\begin{bankq}[id=TOAS04-P2-MATH-S1-Q03, source={TOAS - Test 04 (15 Mar 2022)}, paper=2, qno=3, type=multi, subtopic={Standard limits}, answer={A,B}, answer_src=key, difficulty=4, diff_parts={K1 S2 I1 C0 F1}, diff_note={concepts: limits with greatest integer functions | steps: [x\^{}3] - [x]\^{}3 near each integer; one-sided behaviour; condition for a finite limit | I: case analysis at integers}, also_in={TOAS - Test 07 (01 May 2022) P1 Q5; TEAS - Test 01 (16 Apr 2023) P1 Q2}]
Given $\displaystyle\lim_{x\to\alpha}\dfrac{\sin\left(\left[x^3\right]-[x]^3\right)}{(x-\alpha)^3}$ exists finitely, then the value of $\alpha$ can be (where $\alpha\in I$ \& $[.]$ denotes greatest integer function)-
\begin{qoptions}
\opt{A}{0}
\opt{B}{1}
\opt{C}{$-1$}
\opt{D}{2}
\end{qoptions}
\end{bankq}
\begin{bankq}[id=TOAS04-P2-MATH-S4-Q08, source={TOAS - Test 04 (15 Mar 2022)}, paper=2, qno=8, type=numerical, subtopic={Expansions}, answer={8}, answer_src=key, difficulty=4, diff_parts={K1 S2 I1 C0 F1}, diff_note={concepts: limit of the form 1\^{}infinity with expansion | steps: base tends to 1 so coefficients up to x\^{}3 are fixed; a4 from the exponent 2/x; sum of coefficients | I: match the base to 1 + a4 x/3}]
Let $f(x)=a_0+a_1x+a_2x^2+a_3x^3+a_4x^4$ where $a_i\in R$ and $\displaystyle\lim_{x\to0}\left(\dfrac{f(x)-x+2x^2}{3x^3}\right)^{2/x}=e^4$, then $\displaystyle\sum_{i=0}^{4}a_i$ is equal to
\end{bankq}
\begin{bankq}[id=TOAS05-P1-MATH-S1-Q04, source={TOAS - Test 05 (03 Apr 2022)}, paper=1, qno=4, type=multi, subtopic={Standard limits}, answer={A}, answer_src=key, difficulty=6, diff_parts={K2 S2 I2 C0 F1}, diff_note={concepts: Riemann sums, definite integral | steps: substitute t = e\^{}(-h); the sum becomes a Riemann sum; integral of 1/(e\^{}x + 1) equals ln 2 | I: sum as an integral}]
The value of $\displaystyle\lim_{t\to1^-}(1-t)\sum_{r=1}^{\infty}\dfrac{t^r}{1+t^r}$ is
\begin{qoptions}
\opt{A}{$\ln2$}
\opt{B}{$\ln3$}
\opt{C}{$\ln\left(\dfrac{2e}{1+e}\right)$}
\opt{D}{$\dfrac12$}
\end{qoptions}
\end{bankq}
\begin{bankq}[id=TOAS05-P1-MATH-S3-Q03, source={TOAS - Test 05 (03 Apr 2022)}, paper=1, qno=3, type=numerical, subtopic={Expansions}, answer={2.00}, answer_src=key, difficulty=4, diff_parts={K1 S2 I1 C0 F1}, diff_note={concepts: finite limit forces coefficients | steps: [x/pi] = 0 gives the first limit pi\^{}2; expansion of sin(sin x) - sin x; c = 0; solve b | I: finite limit means c = 0}, also_in={TEAS - Test 01 (16 Apr 2023) P1 Q7}]
If $\displaystyle\lim_{x\to0^+}\left(\dfrac{\sin^2x}{\left[\frac x\pi\right]+\frac{x^2}{\pi^2}}\right)\times\dfrac{\sin(\sin x)-\sin x}{ax^5+bx^3+c}=\dfrac{-\pi^2}{12}$. Then value of $b+c$ is\
(Where $[.]$ denotes greatest integer function)
\end{bankq}
\begin{bankq}[id=TOAS05-P1-MATH-S3-Q04, source={TOAS - Test 05 (03 Apr 2022)}, paper=1, qno=4, type=numerical, subtopic={Standard limits}, answer={0.25}, answer_src=key, difficulty=3, diff_parts={K1 S2 I0 C0 F1}, diff_note={concepts: standard limits | steps: expand cos(pi/x) and the logarithm; tan of a small argument; limit gives -pi\^{}2/4}, also_in={TEAS - Test 01 (16 Apr 2023) P1 Q8}]
If the value of $\displaystyle\lim_{x\to\infty}x^2\tan\left(\ln\left(\sqrt{\cos\dfrac\pi x}\right)\right)$ is equal to $-k\pi^2$, then $k$ is equal to
\end{bankq}
\begin{bankq}[id=TOAS05-P1-MATH-S3-Q05, source={TOAS - Test 05 (03 Apr 2022)}, paper=1, qno=5, type=numerical, subtopic={Standard limits}, answer={2100.00}, answer_src=key, difficulty=4, diff_parts={K1 S2 I1 C0 F1}, diff_note={concepts: limit of a polynomial in tan | steps: set t = tan theta; expand each term; cancel 20 t; sum 10 k | I: leading terms cancel}]
Let \[\begin{aligned}f(\theta)=\dfrac1{\tan^9\theta}\Big(&(1+\tan\theta)^{10}+(2+\tan\theta)^{10}\\&+\ldots.(20+\tan\theta)^{10}\Big)-20\tan\theta.\end{aligned}\] The value of $\displaystyle\lim_{\theta\to\frac\pi2^-}f(\theta)$ is
\end{bankq}
\begin{bankq}[id=TOAS05-P2-MATH-S1-Q08, source={TOAS - Test 05 (03 Apr 2022)}, paper=2, qno=8, type=single, subtopic={Standard limits}, answer={D}, answer_src=key, difficulty=3, diff_parts={K1 S2 I1 C0 F0}, diff_note={concepts: limit with a greatest integer function | steps: [tan x/x] = 1 near 0; expansions of ln(1 + 3x\^{}2), sin(tan\^{}-1 x); (3\^{}x - 1) and (1 - cos 2x) | I: greatest integer equals 1}]
Value of $\displaystyle\lim_{x\to0}\dfrac{\ln\left(1+3x^2\right)\sin\left(\tan^{-1}x\right)}{(3^x-1)(1-\cos2x)}\left[\dfrac{\tan x}{x}\right]$ is (where $[.]$ represents greatest integer function) -
\begin{qoptions}
\opt{A}{$\dfrac32$}
\opt{B}{$\dfrac12$}
\opt{C}{$\dfrac3{\ln3}$}
\opt{D}{$\dfrac3{2\ln3}$}
\end{qoptions}
\end{bankq}
\begin{bankq}[id=TOAS05-P2-MATH-S1-Q09, source={TOAS - Test 05 (03 Apr 2022)}, paper=2, qno=9, type=single, subtopic={Expansions}, answer={D}, answer_src=key, difficulty=2, diff_parts={K1 S2 I0 C0 F0}, diff_note={concepts: finite limit forces constants | steps: expand numerator and denominator; kill x term; x\^{}3 coefficient; solve}]
If $\displaystyle\lim_{x\to0}\dfrac{\alpha\sin x-\beta x\cos x}{\tan x-\sin x}=1$, then $(\alpha+\beta)$ is equal to-
\begin{qoptions}
\opt{A}{1}
\opt{B}{$-1$}
\opt{C}{2}
\opt{D}{3}
\end{qoptions}
\end{bankq}
\begin{bankq}[id=TOAS05-P2-MATH-S1-Q10, source={TOAS - Test 05 (03 Apr 2022)}, paper=2, qno=10, type=single, subtopic={Expansions}, answer={A}, answer_src=key, difficulty=3, diff_parts={K1 S2 I0 C1 F0}, diff_note={concepts: limit by expansions | steps: sin 2x and tan x - tan\^{}-1 x; e\^{}(cos 3x) - e; ln cos x; combine}]
Value of $\displaystyle\lim_{x\to0}\dfrac{(\sin2x)\left(\tan x-\tan^{-1}x\right)}{\left(e^{\cos3x}-e\right)(\ln\cos x)}$ is equal to-
\begin{qoptions}
\opt{A}{$\dfrac{16}{27e}$}
\opt{B}{$\dfrac{16}{27}$}
\opt{C}{$\dfrac4{9e}$}
\opt{D}{$\dfrac49$}
\end{qoptions}
\end{bankq}
\begin{bankq}[id=TOAS05-P2-MATH-S1-Q11, source={TOAS - Test 05 (03 Apr 2022)}, paper=2, qno=11, type=single, subtopic={Expansions}, answer={B}, answer_src=key, difficulty=2, diff_parts={K1 S2 I0 C0 F0}, diff_note={concepts: limit by expansions | steps: expand sqrt(cos 2x); expand sqrt(1 + x sin x); denominator; ratio}]
Value of $\displaystyle\lim_{x\to0}\dfrac{\sqrt{\cos2x}-\sqrt{1+x\sin x}}{x\tan\frac x2}$ is -
\begin{qoptions}
\opt{A}{3}
\opt{B}{$-3$}
\opt{C}{6}
\opt{D}{$-6$}
\end{qoptions}
\end{bankq}
\begin{bankq}[id=TOAS05-P2-MATH-S1-Q12, source={TOAS - Test 05 (03 Apr 2022)}, paper=2, qno=12, type=single, subtopic={Standard limits}, answer={A}, answer_src=key, difficulty=2, diff_parts={K1 S1 I1 C0 F0}, diff_note={concepts: limit with a common factor | steps: factor x - 1 from the denominator; cancel and evaluate | I: factor out x - 1}]
The value of $\displaystyle\lim_{x\to1}\dfrac{x^3-1}{(x^2-3x+2)-\sin(x-1)}$ is-
\begin{qoptions}
\opt{A}{$-\dfrac32$}
\opt{B}{1}
\opt{C}{2}
\opt{D}{3}
\end{qoptions}
\end{bankq}
\begin{bankq}[id=TOAS05-P2-MATH-S1-Q13, source={TOAS - Test 05 (03 Apr 2022)}, paper=2, qno=13, type=single, subtopic={Standard limits}, answer={C}, answer_src=key, difficulty=2, diff_parts={K1 S1 I1 C0 F0}, diff_note={concepts: limit with a periodic function | steps: sec x = 1 + x\^{}2/2; sin(2 pi + pi x\^{}2) ~ pi x\^{}2 | I: reduce modulo 2 pi}]
Value of $\displaystyle\lim_{x\to0}\dfrac{\sin(2\pi\sec x)}{x^2}$ is-
\begin{qoptions}
\opt{A}{$\dfrac\pi4$}
\opt{B}{$\dfrac\pi2$}
\opt{C}{$\pi$}
\opt{D}{$2\pi$}
\end{qoptions}
\end{bankq}
\begin{bankq}[id=TOAS05-P2-MATH-S2-Q06, source={TOAS - Test 05 (03 Apr 2022)}, paper=2, qno=6, type=numerical, subtopic={Expansions}, answer={6}, answer_src=solved, working={$x-\tan x\approx-x^3/3$, $x-\sin x\approx x^3/6$, $x^2-(\sin^{-1}x)^2\approx-x^4/3$: limit $=(-2)\cdot x^2\cdot x^2/(-x^4/3)=6$}, difficulty=4, diff_parts={K1 S2 I0 C1 F1}, diff_note={concepts: limit by expansions | steps: x - tan x and x - sin x; (sin\^{}-1 x)\^{}2 and x\^{}2 - (sin\^{}-1 x)\^{}2; combine}]
Value of $\displaystyle\lim_{x\to0}\dfrac{x^2(x-\tan x)}{(x-\sin x)}\cdot\dfrac{\left(\sin^{-1}x\right)^2}{x^2-\left(\sin^{-1}x\right)^2}$ is
\end{bankq}
\begin{bankq}[id=TOAS05-P2-MATH-S2-Q07, source={TOAS - Test 05 (03 Apr 2022)}, paper=2, qno=7, type=numerical, subtopic={Standard limits}, answer={8}, answer_src=key, difficulty=3, diff_parts={K1 S2 I0 C0 F1}, diff_note={concepts: limits of the form 1\^{}infinity | steps: first limit equals 1/e; second limit equals 1/e; combine}]
The value of $4e\left(\displaystyle\lim_{x\to0}\left(\dfrac{\sin x}{x}\right)^{\frac{\sin x}{x-\sin x}}+\lim_{x\to1}(x)^{\frac1{1-x}}\right)$ is
\end{bankq}
\begin{bankq}[id=TOAS06-P1-MATH-S1-Q02, source={TOAS - Test 06 (17 Apr 2022)}, paper=1, qno=2, type=single, subtopic={Standard limits}, answer={B}, answer_src=key, difficulty=3, diff_parts={K1 S2 I1 C0 F0}, diff_note={concepts: limit of a difference of fractions | steps: substitute x = 1 + h; expand each term; combine to -1/2 | I: expand around 1}]
The value of $\displaystyle\lim_{x\to1}\left(\frac{2018}{1-x^{2018}}-\frac{2019}{1-x^{2019}}\right)$ is
\begin{qoptions}
\opt{A}{0}
\opt{B}{$-\dfrac12$}
\opt{C}{$\dfrac12$}
\opt{D}{does not exist}
\end{qoptions}
\end{bankq}
\begin{bankq}[id=TOAS06-P1-MATH-S2-Q10, source={TOAS - Test 06 (17 Apr 2022)}, paper=1, qno=10, type=multi, subtopic={Standard limits}, answer={B,C,D}, answer_src=key, difficulty=4, diff_parts={K1 S2 I1 C0 F1}, diff_note={concepts: limits with fractional part | steps: 2013\^{}x near 0 from the right; from the left; evaluate x\^{}(1/{x}) in each case | I: 1\^{}infinity and fractional part}]
Let $f(x)=x^{\frac{1}{\{x\}}}$ (where $\{x\}$ denotes fractional part of $x$) then
\begin{qoptions}
\opt{A}{$\displaystyle\lim_{x\to0^+}f(2013^x)=2013$}
\opt{B}{$\displaystyle\lim_{x\to0^-}f(2013^x)=1$}
\opt{C}{$\displaystyle\lim_{x\to0^+}f(2013^x)=e$}
\opt{D}{$\displaystyle\lim_{x\to0}f(2013^x)$ does not exist}
\end{qoptions}
\end{bankq}
\begin{bankpassage}[id=TOAS07-P1-MATH-S3-G1-2, source={TOAS - Test 07 (01 May 2022)}, questions={1-2}]
Given $g(x)=\sin\dfrac1x$, $h(x)=x^2+2x+(\lambda+1)$ and $u(x)=\dfrac1x+\cos\dfrac1{x^2}$\\
Let $f(x)=\displaystyle\lim_{n\to\infty}\dfrac{x^{2n+1}g(x)+h(x)}{x^{2n}+3x.\,u(x)}$
\end{bankpassage}
\begin{bankq}[id=TOAS07-P1-MATH-S3-Q01, source={TOAS - Test 07 (01 May 2022)}, paper=1, qno=1, type=numerical, subtopic={Standard limits}, answer={5.00}, answer_src=key, difficulty=3, diff_parts={K1 S2 I0 C0 F1}, diff_note={concepts: limit problem from the passage | steps: identify f; expansions near 0; value of lambda}, group=TOAS07-P1-MATH-S3-G1-2]
If $\displaystyle\lim_{x\to0}f(x)=2$, then the value of $\lambda$ is -
\end{bankq}
\begin{bankq}[id=TOAS07-P1-MATH-S3-Q02, source={TOAS - Test 07 (01 May 2022)}, paper=1, qno=2, type=numerical, subtopic={Standard limits}, answer={1.00}, answer_src=key, difficulty=2, diff_parts={K1 S1 I0 C0 F1}, diff_note={concepts: limit at infinity from the passage | steps: leading terms of f; limit}, group=TOAS07-P1-MATH-S3-G1-2]
$\displaystyle\lim_{x\to\infty}f(x)$ is equal to -
\end{bankq}
\begin{bankq}[id=TOAS07-P2-MATH-S1-Q05, source={TOAS - Test 07 (01 May 2022)}, paper=2, qno=5, type=multi, subtopic={Standard limits}, answer={A,B,D}, answer_src=key, difficulty=4, diff_parts={K1 S2 I1 C0 F1}, diff_note={concepts: limit of sin\^{}-1 sin | steps: argument tends to 1 from above; reduce by the principal branch; compare with the options | I: principal branch of sin\^{}-1 sin}, also_in={TNAS - Test 02 (07 May 2023) P2 Q6}]
$\displaystyle\lim_{x\to\infty}\sin^{-1}\sin\left(\dfrac{x^2+2}{x^2+1}\right)$ is equal to (where $[.]$ and $\{.\}$ denote greatest integer and fractional part function respectively)
\begin{qoptions}
\opt{A}{$\displaystyle\lim_{x\to\infty}\left[\dfrac{x^2+2}{x^2+1}\right]$}
\opt{B}{$\displaystyle\lim_{x\to\infty}\dfrac{x^2+2}{x^2+1}$}
\opt{C}{$\displaystyle\lim_{x\to\infty}\left\{\dfrac{x^2+2}{x^2+1}\right\}$}
\opt{D}{$\displaystyle\lim_{x\to\infty}\left\{\dfrac{x^2}{x^2+1}\right\}$}
\end{qoptions}
\end{bankq}
\begin{bankpassage}[id=TOAS07-P2-MATH-S3-G5-6, source={TOAS - Test 07 (01 May 2022)}, questions={5-6}]
Consider $f(x)=x^4-3x^3+2$ and let $h(x)=\dfrac{f(x)}{g(x)}$ be a rational function such that\\
(i) $h(x)$ is continuous in $R-\{-1,1\}$.\\
(ii) $g(x)$ does not have any repeated linear factors.\\
(iii) As $x$ approaches infinity, $h(x)$ also approaches infinity.\\
(iv) $\displaystyle\lim_{x\to1}h(x)=\dfrac12$
\end{bankpassage}
\begin{bankq}[id=TOAS07-P2-MATH-S3-Q05, source={TOAS - Test 07 (01 May 2022)}, paper=2, qno=5, type=numerical, subtopic={Standard limits}, answer={5.00}, answer_src=key, difficulty=2, diff_parts={K1 S1 I0 C0 F1}, diff_note={concepts: range of a function from the passage | steps: express g; maximum value}, group=TOAS07-P2-MATH-S3-G5-6]
Range of $g(x)$ is $(-\infty,k]$ then $k$ is equal to
\end{bankq}
\begin{bankq}[id=TOAS07-P2-MATH-S3-Q06, source={TOAS - Test 07 (01 May 2022)}, paper=2, qno=6, type=numerical, subtopic={Standard limits}, answer={1.00}, answer_src=key, difficulty=3, diff_parts={K1 S2 I0 C0 F1}, diff_note={concepts: limit from the passage | steps: substitute g(cos x); expand near 0; ratio}, group=TOAS07-P2-MATH-S3-G5-6]
$\displaystyle\lim_{x\to0}\dfrac{g(\cos x)}{5\ln\left(\sec^2x\right)}$ is equal to -
\end{bankq}
\begin{bankq}[id=TOAS08-P1-MATH-S1-Q03, source={TOAS - Test 08 (15 May 2022)}, paper=1, qno=3, type=multi, subtopic={Expansions}, answer={A,D}, answer_src=key, difficulty=4, diff_parts={K1 S2 I1 C0 F1}, diff_note={concepts: limit of (sec a)\^{}(1/x) differences | steps: expand (sec a)\^{}(1/x) near a = 0; limit gives f(x,y) = 2xy/(y - x); examine the curve f = 12 | I: expand to a\^{}2}, also_in={TMAS - Test 02 (06 Aug 2023) P1 Q5; TNAS - Test 04 (18 Jun 2023) P2 Q1}]
Let $f(x,y)=\dfrac{1}{\displaystyle\lim_{\alpha\to0}\left(\dfrac{\sqrt[x]{\sec\alpha}-\sqrt[y]{\sec\alpha}}{\alpha^2}\right)}$, then
\begin{qoptions}
\opt{A}{$f(2,3)=12$}
\opt{B}{$f(x,y)=12$ has infinitely many solutions}
\opt{C}{Slope of tangent to the curve $f(x,y)=12$ at $x=2$ is 4}
\opt{D}{$f(x,y)=12$ has no horizontal tangent}
\end{qoptions}
\end{bankq}
\begin{bankq}[id=TOAS08-P1-MATH-S3-Q04, source={TOAS - Test 08 (15 May 2022)}, paper=1, qno=4, type=numerical, subtopic={Expansions}, answer={8}, answer_src=solved, working={$\sqrt{x^2+\sqrt{1+x^4}}=\sqrt2x(1+\frac1{8x^4}+\dots)$ so the limit is $\sqrt2/8$, $\lambda=8$}, difficulty=5, diff_parts={K1 S2 I1 C1 F1}, diff_note={concepts: limit at infinity with nested radicals | steps: factor sqrt2 x; expand sqrt(1 + 1/x\^{}4); expand the outer radical; leading term | I: factor out the dominant term}, also_in={TMAS - Test 02 (06 Aug 2023) P1 Q7; TNAS - Test 04 (18 Jun 2023) P2 Q4}]
$\displaystyle\lim_{x\to\infty}x^3\left(\sqrt{x^2+\sqrt{1+x^4}}-\sqrt2x\right)=\dfrac{\sqrt2}{\lambda}$ find value of $\lambda$
\end{bankq}
\begin{bankq}[id=TOAS08-P2-MATH-S2-Q08, source={TOAS - Test 08 (15 May 2022)}, paper=2, qno=8, type=single, subtopic={Standard limits}, answer={C}, answer_src=key, difficulty=2, diff_parts={K1 S1 I0 C1 F0}, diff_note={concepts: limits matching | steps: evaluate the four limits; match with List-II}]
Match List-I with List-II and select the correct answer using the code given below the list.
\begin{center}\footnotesize\renewcommand{\arraystretch}{1.1}\setlength{\extrarowheight}{3.5pt}\setlength{\tabcolsep}{2pt}
\begin{tabular}{|c|l|c|c|}
\hline
 & \multicolumn{1}{c|}{List-I} & & List-II\\[3pt]
\hline
(P) & \parbox[t]{4.4cm}{If $x=\alpha$ satisfies $\sin^{-1}x+\cos^{-1}x^2+\dfrac\pi2=0$, then the value of $\dfrac{\sec^{-1}\alpha-\tan^{-1}\alpha}{\cot^{-1}\alpha-\mathrm{cosec}^{-1}\alpha}$ is} & (1) & 0\\[14pt]
(Q) & \parbox[t]{4.4cm}{If $a=\displaystyle\lim_{x\to\infty}\left(2.3^x+3.4^x+4.5^x\right)^{\frac1x}$ and $b=\displaystyle\lim_{x\to(-\infty)}\left(2.3^x+3.4^x+4.5^x\right)^{\frac1x}$ then $\left(\dfrac{a+b}{a-b}\right)$ is equal to} & (2) & 4\\[22pt]
(R) & \parbox[t]{4.4cm}{Number of points of discontinuity of $f(x)=x\sin(\pi(2x-[x]))$ in $[0,2]$ will be (Note $[k]$ denotes greatest integer function less than or equal to $k$)} & (3) & 1\\[22pt]
(S) & \parbox[t]{4.4cm}{$\displaystyle\lim_{x\to0}\left(\dfrac{3e^{x/3}-x-3}{x^2}\right)$ is equal to} & (4) & $\dfrac13$\\[6pt]
 & & (5) & $\dfrac12$\\[6pt]
 & & (6) & $\dfrac16$\\[3pt]
\hline
\end{tabular}
\end{center}
\begin{qoptions}
\opt{A}{$P\to1;Q\to2;R\to2;S\to4$}
\opt{B}{$P\to3;Q\to1;R\to2;S\to6$}
\opt{C}{$P\to3;Q\to2;R\to1;S\to6$}
\opt{D}{$P\to1;Q\to3;R\to2;S\to4$}
\end{qoptions}
\end{bankq}
\begin{bankq}[id=TOAS08-P2-MATH-S3-Q04, source={TOAS - Test 08 (15 May 2022)}, paper=2, qno=4, type=numerical, subtopic={Standard limits}, answer={0.50}, answer_src=key, difficulty=3, diff_parts={K1 S1 I1 C0 F1}, diff_note={concepts: telescoping infinite product | steps: 1 - 1/k\^{}2 = (k - 1)(k + 1)/k\^{}2; product telescopes to (n + 1)/(2n) | I: factor each term}]
The value of\[\lim_{n\to\infty}\left(1-\dfrac1{2^2}\right)\left(1-\dfrac1{3^2}\right)\left(1-\dfrac1{4^2}\right)\ldots\left(1-\dfrac1{n^2}\right)\]is
\end{bankq}
\begin{bankq}[id=TOAS09-P2-MATH-S1-Q08, source={TOAS - Test 09 (30 May 2022)}, paper=2, qno=8, type=single, subtopic={Standard limits}, answer={B}, answer_src=key, difficulty=1, diff_parts={K1 S1 I0 C0 F0}, diff_note={concepts: limit with small-angle approximations | steps: numerator ~ 6x, denominator ~ 2x}, also_in={TNAS - Test 05 (09 Jul 2023) P1 Q8}]
$\displaystyle\lim_{x\to0}\dfrac{5\sin x+x\cos x}{2\tan x-x^2}$ is
\begin{qoptions}
\opt{A}{$\dfrac52$}
\opt{B}{3}
\opt{C}{0}
\opt{D}{does not exist}
\end{qoptions}
\end{bankq}
\begin{bankq}[id=TOAS09-P2-MATH-S1-Q15, source={TOAS - Test 09 (30 May 2022)}, paper=2, qno=15, type=single, subtopic={Standard limits}, answer={A}, answer_src=key, difficulty=3, diff_parts={K1 S2 I1 C0 F0}, diff_note={concepts: limit with binomial expansion | steps: factor x\^{}(2/3) out; binomial expansion of (1 + 1/x)\^{}(2/3) - (1 - 1/x)\^{}(2/3); limit 4/3 | I: binomial expansion at infinity}, also_in={TNAS - Test 05 (09 Jul 2023) P1 Q15}]
The value of $\displaystyle\lim_{x\to\infty}x^{\frac13}\left((x+1)^{\frac23}-(x-1)^{\frac23}\right)$ is
\begin{qoptions}
\opt{A}{$\dfrac43$}
\opt{B}{$\dfrac{-1}{3}$}
\opt{C}{$\dfrac13$}
\opt{D}{$\dfrac{-2}{3}$}
\end{qoptions}
\end{bankq}
\begin{bankq}[id=TOAS09-P2-MATH-S2-Q03, source={TOAS - Test 09 (30 May 2022)}, paper=2, qno=3, type=numerical, subtopic={Standard limits}, answer={4}, answer_src=solved, working={Both arguments tend to 1 and 2 where $f=8$, 12: ratio $=\frac{f'(1)\cdot4}{f'(2)\cdot3}=\frac{6\cdot4}{2\cdot3}=4$}, difficulty=4, diff_parts={K1 S2 I1 C0 F1}, diff_note={concepts: limit as a ratio of derivatives | steps: arguments tend to 1 and 2; f(1) = 8 and f(2) = 12; use derivatives of the arguments | I: values at limit points}]
If $f(x)=10x-2x^2$, then
\[\lim_{h\to0}\frac{f(3\cosh h+4\sinh h-2)-8}{f(3e^h-5\sec h+4)-12}\]
is equal to
\end{bankq}
\begin{bankq}[id=TOAS09-P2-MATH-S2-Q04, source={TOAS - Test 09 (30 May 2022)}, paper=2, qno=4, type=numerical, subtopic={Expansions}, answer={512}, answer_src=solved, working={Continuity needs $b=-1$, $a=-4$ (the $x$ and $x^2$ terms vanish); $L=a^3/6+64/3=32/3$; $12abL=512$}, difficulty=5, diff_parts={K1 S3 I0 C1 F1}, diff_note={concepts: continuity and series expansion | steps: expand e\^{}(ax) and ln(1 + 4x); kill the x term; kill the x\^{}2 term; L from x\^{}3; 12abL}, also_in={TMAS - Test 01 (09 Jul 2023) P1 Q2; TNAS - Test 03 (28 May 2023) P1 Q2}]
If $f(x)=\begin{cases}\dfrac{e^{ax}+\ell n(1+4x)+b}{x^3} & x\neq0\\ L & x=0\end{cases}$ is continuous at $x=0$, then $12abL$ is equal to
\end{bankq}
\begin{bankq}[id=TOAS10-P1-MATH-S1-Q01, source={TOAS - Test 10 (12 Jun 2022)}, paper=1, qno=1, type=multi, subtopic={Standard limits}, answer={C}, answer_src=key, difficulty=2, diff_parts={K1 S1 I0 C0 F1}, diff_note={concepts: standard limits | steps: expand cos(pi/x) and the logarithm; tan of the small argument}]
$\displaystyle\lim_{x\to\infty} x^2\tan\left(\ln\sqrt{\cos\dfrac{\pi}{x}}\right)$ is equal to
\begin{qoptions}
\opt{A}{$-\dfrac{\pi^2}{2}$}
\opt{B}{$\dfrac{\pi^2}{2}$}
\opt{C}{$-\dfrac{\pi^2}{4}$}
\opt{D}{$\dfrac{\pi^2}{4}$}
\end{qoptions}
\end{bankq}
\begin{bankq}[id=TOAS10-P1-MATH-S3-Q06, source={TOAS - Test 10 (12 Jun 2022)}, paper=1, qno=6, type=numerical, subtopic={Standard limits}, answer={1}, answer_src=solved, working={With s = sin theta the exponent is s(1 - s) and s\^{}(s(1-s)) = exp(s(1-s) ln s), which tends to exp(0) = 1}, difficulty=4, diff_parts={K1 S2 I1 C0 F1}, diff_note={concepts: limit of the form s to the power s(1 - s) | steps: set s = sin theta; s\^{}(s(1-s)) = exp(s(1-s) ln s); limit is 1 | I: s ln s tends to 0}]
Find the value of $\displaystyle\lim_{\theta\to0^+}(\sin\theta)^{\sin\theta-\sin^2\theta}$.
\end{bankq}
\begin{bankq}[id=TOAS11-P1-MATH-S2-Q11, source={TOAS - Test 11 (27 Jun 2022)}, paper=1, qno=11, type=multi, subtopic={Standard limits}, answer={A,B,C}, answer_src=key, difficulty=4, diff_parts={K1 S2 I1 C0 F1}, diff_note={concepts: limit of a polynomial in tan | steps: set t = tan theta; expand each term; cancel 20 t; sum 10 k | I: leading terms cancel}, also_in={TNAS - Test 05 (09 Jul 2023) P2 Q7}]
Let
\[\begin{aligned}f(\theta)=\dfrac{1}{\tan^9\theta}&\left((1+\tan\theta)^{10}+(2+\tan\theta)^{10}\right.\\&\left.+\ldots+(20+\tan\theta)^{10}\right)-20\tan\theta.\end{aligned}\]
If $\displaystyle\lim_{\theta\to\frac\pi2^-}f(\theta)$ is $\ell$, then
\begin{qoptions}
\opt{A}{$\ell$ is divisible by 100}
\opt{B}{$\ell$ is divisible by 7}
\opt{C}{$\ell$ is divisible by 30}
\opt{D}{$\ell$ is a prime number}
\end{qoptions}
\end{bankq}
\begin{bankq}[id=TOAS12-P1-MATH-S2-Q04, source={TOAS - Test 12 (10 Jul 2022)}, paper=1, qno=4, type=numerical, subtopic={L'Hopital}, answer={6.00}, answer_src=key, difficulty=3, diff_parts={K1 S2 I0 C0 F1}, diff_note={concepts: limit with Taylor expansion of f | steps: expand f(5x), f(3x), f(2x); constant and linear terms cancel; coefficient of x\^{}2}]
Let $f(x)$ be a twice differentiable function such that $f''(0)=2$, then the value of $\displaystyle\lim_{x\to0}\frac{f(5x)-3f(3x)+2f(2x)}{x^2}$ is
\end{bankq}
\begin{bankpassage}[id=TOAS12-P2-MATH-S2-G9-10, source={TOAS - Test 12 (10 Jul 2022)}, questions={9-10}]
Let $f(x)=ax^3+bx^2+cx+d$, satisfying $\displaystyle\lim_{x\to0}\left(3+\dfrac{f(x)}{x^2}\right)^{\frac1x}=e^2$, then\par
On the basis of above information, answer the following questions :
\end{bankpassage}
\begin{bankq}[id=TOAS12-P2-MATH-S2-Q09, source={TOAS - Test 12 (10 Jul 2022)}, paper=2, qno=9, type=multi, subtopic={Standard limits}, answer={A,B,C,D}, answer_src=key, difficulty=3, diff_parts={K1 S1 I1 C0 F1}, diff_note={concepts: limit of the form 1\^{}infinity for a cubic | steps: exponent condition gives a, b relations; evaluate the options | I: expand the base to find coefficients}, group=TOAS12-P2-MATH-S2-G9-10]
Correct option's is /are -
\begin{qoptions}
\opt{A}{$a+b=0$}
\opt{B}{$a+b+c+d=0$}
\opt{C}{$a+c=2$}
\opt{D}{$d-b=2$}
\end{qoptions}
\end{bankq}
\begin{bankq}[id=TOAS12-P2-MATH-S2-Q10, source={TOAS - Test 12 (10 Jul 2022)}, paper=2, qno=10, type=multi, subtopic={Standard limits}, answer={A}, answer_src=key, difficulty=3, diff_parts={K1 S2 I0 C0 F1}, diff_note={concepts: sum of roots for several equations | steps: identify f; sum of roots of each cubic from Vieta; sum over r}, group=TOAS12-P2-MATH-S2-G9-10]
Sum of all roots of the all possible equations $f(x)-2(r^2-1)x^2=0$ (where $r=1,2,\dots,10$)
\begin{qoptions}
\opt{A}{385}
\opt{B}{250}
\opt{C}{165}
\opt{D}{305}
\end{qoptions}
\end{bankq}
\begin{bankpassage}[id=TOAS13-P1-MATH-S2-G9-10, source={TOAS - Test 13 (24 Jul 2022)}, questions={9-10}]
Suppose that function $f:R\to R$ satisfies the inequality, $\left|\displaystyle\sum_{k=1}^{n}\left(3^k\{f(x+ky)-f(x-ky)\}\right)\right|\le1$ for every positive integer $n$ and for all $x,y\in R$. Also, let $g(x)=ax^3-(3+b)x^2+cx$ for $x\in R$, where $a$, $b$, $c$ are such that $\displaystyle\lim_{x\to0}\dfrac{axe^x-b\ln(1+x)+cxe^{-x}}{x^2\sin x}=2$
\end{bankpassage}
\begin{bankq}[id=TOAS13-P1-MATH-S2-Q09, source={TOAS - Test 13 (24 Jul 2022)}, paper=1, qno=9, type=multi, subtopic={Expansions}, answer={A,B,D}, answer_src=key, difficulty=8, diff_parts={K3 S3 I1 C1 F1}, diff_note={concepts: telescoping inequality, series expansions in a limit, counting roots of a cubic | steps: f is constant from the n and n-1 sums and scaling y; expand the limit to order x\^{}3; three linear equations give a = 3, b = 12, c = 9; critical points of g; values at the critical points against -23; count solutions}, group=TOAS13-P1-MATH-S2-G9-10]
If $f(13)=-23$, number of solutions of $g(x)=f(x)$ is more than or equal to :
\begin{qoptions}
\opt{A}{2}
\opt{B}{3}
\opt{C}{4}
\opt{D}{1}
\end{qoptions}
\end{bankq}
\begin{bankq}[id=TOAS13-P1-MATH-S2-Q10, source={TOAS - Test 13 (24 Jul 2022)}, paper=1, qno=10, type=multi, subtopic={Expansions}, answer={A}, answer_src=key, difficulty=5, diff_parts={K2 S2 I0 C1 F1}, diff_note={concepts: Taylor expansion of the limit, L'Hopital rule | steps: a, b, c from the expansion; g(1) = -3 so the limit is 0/0; differentiate numerator and denominator at x = 1}, group=TOAS13-P1-MATH-S2-G9-10]
The value of $\displaystyle\lim_{x\to1}\dfrac{g(x)+3}{\ln(2-x)}$ is equal to
\begin{qoptions}
\opt{A}{12}
\opt{B}{3}
\opt{C}{9}
\opt{D}{1}
\end{qoptions}
\end{bankq}
\begin{bankq}[id=TOAS15-P1-MATH-S1-Q05, source={TOAS - Test 15 (22 Aug 2022)}, paper=1, qno=5, type=single, subtopic={Standard limits}, answer={B}, answer_src=key, difficulty=4, diff_parts={K2 S2 I1 C0 F0}, diff_note={concepts: limit with the greatest integer function, expansion of the square root | steps: sqrt(n\^{}2 + n + 1) lies between n and n+1; floor is n; limit of the fractional part is 1/2}]
$\displaystyle\lim_{n\to\infty}\left(\sqrt{n^2+n+1}-\left[\sqrt{n^2+n+1}\right]\right)$ $(n\in N)$ (where $[.]$ denotes the greatest integer function) is
\begin{qoptions}
\opt{A}{0}
\opt{B}{1/2}
\opt{C}{2/3}
\opt{D}{1/4}
\end{qoptions}
\end{bankq}
\begin{bankpassage}[id=TOAS16-P2-MATH-S2-G9-10, source={TOAS - Test 16 (04 Sep 2022)}, questions={9-10}]
Consider $f(x)=x^{\ell nx}$ and $g(x)=e^2x$. Let $\alpha$ and $\beta$ ($\alpha<\beta$) be two values of $x$ satisfying equation $f(x)=g(x)$.
\end{bankpassage}
\begin{bankq}[id=TOAS16-P2-MATH-S2-Q09, source={TOAS - Test 16 (04 Sep 2022)}, paper=2, qno=9, type=single, subtopic={Standard limits}, answer={C}, answer_src=key, difficulty=3, diff_parts={K1 S2 I1 C0 F0}, diff_note={concepts: logarithmic equation | steps: take logs; quadratic in ln x; product of roots}, group=TOAS16-P2-MATH-S2-G9-10]
Product of $\alpha\beta$ is -
\begin{qoptions}
\opt{A}{$\dfrac1{e^2}$}
\opt{B}{$\dfrac1e$}
\opt{C}{$e$}
\opt{D}{$e^2$}
\end{qoptions}
\end{bankq}
\begin{bankq}[id=TOAS16-P2-MATH-S2-Q10, source={TOAS - Test 16 (04 Sep 2022)}, paper=2, qno=10, type=single, subtopic={Standard limits}, answer={B}, answer_src=key, difficulty=5, diff_parts={K2 S2 I1 C1 F0}, diff_note={concepts: limit, derivative of a power function | steps: c from the vanishing numerator; derivatives of f and g at beta; ratio}, group=TOAS16-P2-MATH-S2-G9-10]
If $\displaystyle\lim_{x\to\beta}\frac{f(x)-c\beta}{g(x)-\beta^2}$ exists and equal to `$r$' then value of $(c-r)$ is -
\begin{qoptions}
\opt{A}{$4-e^2$}
\opt{B}{$e^2-4$}
\opt{C}{$4-e$}
\opt{D}{$e-4$}
\end{qoptions}
\end{bankq}
\begin{bankq}[id=TOAS17-P2-MATH-S1-Q18, source={TOAS - Test 17 (18 Sep 2022)}, paper=2, qno=18, type=single, subtopic={Standard limits}, answer={A}, answer_src=key, difficulty=3, diff_parts={K1 S1 I1 C1 F0}, diff_note={concepts: limits of exponent forms | steps: each term at zero; combine}]
$\displaystyle\lim_{x\to0^+}\left((\sin x)^{\frac1x}+(\cos x)^{\frac1x}+(\operatorname{cosec}x)^x\right)$ is equal to
\begin{qoptions}
\opt{A}{2}
\opt{B}{4}
\opt{C}{6}
\opt{D}{8}
\end{qoptions}
\end{bankq}
\begin{bankq}[id=TOAS18-P1-MATH-S3-Q04, source={TOAS - Test 18 (30 Oct 2022)}, paper=1, qno=4, type=numerical, subtopic={Standard limits}, answer={1}, answer_src=solved, working={The sum tends to $\int_0^1\frac{dt}{t^2-x^2}=\tfrac1{2x}\ln\tfrac{1-x}{1+x}\to-1$ as $x\to0$; modulus 1}, difficulty=6, diff_parts={K1 S2 I2 C1 F1}, diff_note={concepts: limit of a Riemann sum with a parameter | steps: convert to an integral; principal value; limit as x tends to zero}]
\par\noindent\hfill$\displaystyle\begin{aligned}\Bigl|\lim_{x\to0}\Bigl[\lim_{n\to\infty}\Bigl\{&\dfrac n{1-n^2x^2}+\dfrac n{4-n^2x^2}\\&+\dots+\dfrac n{n^2-n^2x^2}\Bigr\}\Bigr]\Bigr|\end{aligned}$\hfill\null\par
is equal to
\end{bankq}
\begin{bankq}[id=TOAS18-P2-MATH-S1-Q01, source={TOAS - Test 18 (30 Oct 2022)}, paper=2, qno=1, type=single, subtopic={Standard limits}, answer={B}, answer_src=key, difficulty=3, diff_parts={K1 S2 I1 C0 F0}, diff_note={concepts: limit of a product | steps: logarithm of the product; Riemann sum; compare with exponentials}]
Let $f:(0,\infty)\to R$ be a function such that $f(x)=\displaystyle\lim_{n\to\infty}\left(1+\dfrac x{n^2}\right)\left(1+\dfrac{2x}{n^2}\right)\left(1+\dfrac{3x}{n^2}\right)\dots\left(1+\dfrac{nx}{n^2}\right)$. Then correct option is
\begin{qoptions}
\opt{A}{$f'(1)>f(1)$}
\opt{B}{$f(x)<e^x$, $\forall\,x\in(0,\infty)$}
\opt{C}{$f(x)<x^2$, $\forall\,x\in(0,\infty)$}
\opt{D}{$f(c)<\dfrac12$, for some $c$ in $(0,\infty)$}
\end{qoptions}
\end{bankq}

\section*{TAAS papers}
\begin{bankq}[id=TAAS02-P1-MATH-S4-Q06, source={TAAS - Test 02 (30 Jul 2023)}, paper=1, qno=6, type=numerical, subtopic={Standard limits}, answer={0}, answer_src=key, difficulty=6, diff_parts={K1 S2 I2 C1 F1}, diff_note={concepts: limit by substitution and expansion | steps: put t = x - pi/4; expand; evaluate}, also_in={TNAS - Test 06 (30 Jul 2023) P1 Q6}]
The value of
\par\noindent\hfill{\small$\displaystyle\lim_{x\to\frac\pi4}\left(\dfrac{4\sqrt2-(\cos x+\sin x)^5-10\cos x+10\cos x\sin2x}{1-\sin2x}\right)$}\hfill\null\par\noindent is
\end{bankq}
\begin{bankq}[id=TAAS02-P2-MATH-S2-Q06, source={TAAS - Test 02 (30 Jul 2023)}, paper=2, qno=6, type=multi, subtopic={Standard limits}, answer={A,B,C}, answer_src=key, difficulty=6, diff_parts={K1 S2 I2 C1 F1}, diff_note={concepts: limit of a product with exponents | steps: take logarithm; Riemann sum; value; compare with the other limit}, also_in={TNAS - Test 06 (30 Jul 2023) P2 Q6}]
If \par\noindent\hfill$\displaystyle\begin{aligned}p=\lim_{n\to\infty}n^{-n^2}\Bigl\{&(n+2^0)(n+2^{-1})\\&(n+2^{-2})\dots\dots(n+2^{-n+1})\Bigr\}^n\end{aligned}$\hfill\null\par\noindent then which of the following is true ?
\begin{qoptions}
\opt{A}{$p$ is irrational}
\opt{B}{integer nearest to $p$ is 7}
\opt{C}{$p>\displaystyle\lim_{x\to\infty}\left(\dfrac{1+x}x\right)^{\ell nx-x}$}
\opt{D}{None of these}
\end{qoptions}
\end{bankq}

\section*{TEAS papers}
\begin{bankq}[id=TEAS01-P1-MATH-S1-Q01, source={TEAS - Test 01 (16 Apr 2023)}, paper=1, qno=1, type=multi, subtopic={Standard limits}, answer={C}, answer_src=key, difficulty=7, diff_parts={K2 S2 I2 C1 F1}, diff_note={concepts: binomial form, limit | steps: conjugate relation; ratio A over B; limit; compare with the options}]
Let $(1+\sqrt2)^n=A_n+B_n\sqrt2$, where $A_n$ and $B_n$ are rational numbers and $n\in N$, then the value of $\displaystyle\lim_{n\to\infty}\dfrac{A_n}{B_n}$ is equal to -\\
(where [.] denotes greatest integer function)
\begin{qoptions}
\opt{A}{$\left[\displaystyle\lim_{x\to0}\dfrac{\sin x}x\right]$}
\opt{B}{$\left[\displaystyle\lim_{x\to\infty}\dfrac{\sin x}x\right]$}
\opt{C}{$\displaystyle\lim_{x\to\frac\pi4}\dfrac{|\cos x-\sin x|}{|x-\frac\pi4|}$}
\opt{D}{$\left[\displaystyle\lim_{x\to0}\dfrac{\sin\sqrt2x}x\right]$}
\end{qoptions}
\end{bankq}
\begin{bankq}[id=TEAS01-P1-MATH-S3-Q01, source={TEAS - Test 01 (16 Apr 2023)}, paper=1, qno=1, type=numerical, subtopic={Expansions}, answer={8.00}, answer_src=key, difficulty=9, diff_parts={K2 S3 I2 C2 F1}, diff_note={concepts: limit of the form one to infinity, series expansion | steps: exponential limit form; expand to order three; match the coefficients; find the polynomial coefficients; sum}]
Let $f(x)=a_0+a_1x+a_2x^2+a_3x^3+a_4x^4$ where $a_i\in R$ and $\displaystyle\lim_{x\to0}\left(\dfrac{f(x)-x+2x^2}{3x^3}\right)^{\frac2x}=e^4$, then $\displaystyle\sum_{i=0}^{4}a_i$ is equal to
\end{bankq}
\begin{bankq}[id=TEAS01-P1-MATH-S3-Q04, source={TEAS - Test 01 (16 Apr 2023)}, paper=1, qno=4, type=numerical, subtopic={Expansions}, answer={3.00}, answer_src=key, difficulty=7, diff_parts={K1 S3 I2 C1 F1}, diff_note={concepts: limit with expansions | steps: expand log and sine; kill lower order terms; find a and b; limit l; combine}]
If $\displaystyle\lim_{x\to0}\dfrac{a+\log_e(1+x)+b\sin x}{x^2}=\ell$ (finite) then $|a+b+4\ell|$ (where $a,b,\ell\in R$)
\end{bankq}

\section*{TMAS papers}
\begin{bankq}[id=TMAS01-P1-MATH-S1-Q15, source={TMAS - Test 01 (09 Jul 2023)}, paper=1, qno=15, type=single, subtopic={Standard limits}, answer={B}, answer_src=key, difficulty=5, diff_parts={K1 S2 I2 C1 F0}, diff_note={concepts: limit with oscillating term | steps: bound the oscillating term; leading term; the exponent}, also_in={TNAS - Test 03 (28 May 2023) P1 Q16}]
The value of $\displaystyle\lim_{\theta\to0^+}\left(\theta^2\sqrt\theta\cdot\sin\left(\dfrac1\theta\right)+4\sin\left(\theta\sqrt\theta\right)\right)^\theta$ is
\begin{qoptions}
\opt{A}{$-1$}
\opt{B}{1}
\opt{C}{2}
\opt{D}{$e$}
\end{qoptions}
\end{bankq}
\begin{bankq}[id=TMAS01-P1-MATH-S1-Q20, source={TMAS - Test 01 (09 Jul 2023)}, paper=1, qno=20, type=single, subtopic={Standard limits}, answer={C}, answer_src=key, difficulty=5, diff_parts={K1 S2 I2 C1 F0}, diff_note={concepts: limit at minus infinity | steps: rationalise; leading terms; solve}]
$\displaystyle\lim_{x\to-\infty}\left(\sqrt{x^2-x+1}-ax+2b\right)=0$, then $a=$ \underline{\hspace{1.2cm}}.
\begin{qoptions}
\opt{A}{$a=0$}
\opt{B}{$a=\pm2$}
\opt{C}{$a=-1$}
\opt{D}{$a=1$}
\end{qoptions}
\end{bankq}
\begin{bankq}[id=TMAS01-P1-MATH-S2-Q10, source={TMAS - Test 01 (09 Jul 2023)}, paper=1, qno=10, type=numerical, subtopic={Standard limits}, answer={3}, answer_src=key, difficulty=5, diff_parts={K1 S2 I1 C1 F1}, diff_note={concepts: limit with cosine | steps: expansion; standard limit; evaluate}, also_in={TNAS - Test 03 (28 May 2023) P1 Q10}]
If $\displaystyle\lim_{x\to2}\dfrac{1-\cos\left(x^2-6x+8\right)}{(x-2)^2}=p$, then value of $\left(\dfrac6p\right)$ is equal to
\end{bankq}
\begin{bankq}[id=TMAS02-P2-MATH-S1-Q04, source={TMAS - Test 02 (06 Aug 2023)}, paper=2, qno=4, type=multi, subtopic={Standard limits}, answer={B,C}, answer_src=solved, working={$f=\dfrac{x+2}{3-x^2}$ inside $|x|<1$, $2$ outside, $f(1)=5/3$ so it is discontinuous at 1, and the limit at 2 is 2}, difficulty=6, diff_parts={K1 S2 I2 C1 F1}, diff_note={concepts: limit defining a function | steps: cases of x; value at one; continuity; limit at two}]
If $f(x)=\displaystyle\lim_{n\to\infty}\dfrac{2x^{2n}+x+2}{x^{2n}-x^2+3}$, then which of the following statement(s) holds good ?
\begin{qoptions}
\opt{A}{$f(x)$ is continuous at $x=1$}
\opt{B}{$f(x)$ is discontinuous at $x=1$}
\opt{C}{$\displaystyle\lim_{x\to2}f(x)$ exists finitely}
\opt{D}{$\displaystyle\lim_{x\to2}f(x)$ does not exist}
\end{qoptions}
\end{bankq}

\section*{TNAS papers}
\begin{bankq}[id=TNAS02-P1-MATH-S1-Q03, source={TNAS - Test 02 (07 May 2023)}, paper=1, qno=3, type=single, subtopic={Standard limits}, answer={D}, answer_src=key, difficulty=5, diff_parts={K1 S2 I2 C1 F0}, diff_note={concepts: limit with exponents | steps: expand each power to the needed order; subtract; limit}]
The value of {\scriptsize$\displaystyle\lim_{x\to0}\dfrac{(1+3x+2x^2)^{1/x}-(1+3x-2x^2)^{1/x}}{x}$} is
\begin{qoptions}
\opt{A}{$-e^3$}
\opt{B}{$3e^2$}
\opt{C}{$\dfrac13e^{-2}$}
\opt{D}{$4e^3$}
\end{qoptions}
\end{bankq}
\begin{bankq}[id=TNAS02-P1-MATH-S2-Q06, source={TNAS - Test 02 (07 May 2023)}, paper=1, qno=6, type=multi, subtopic={Standard limits}, answer={A,C,D}, answer_src=key, difficulty=6, diff_parts={K1 S2 I2 C1 F1}, diff_note={concepts: limit function with exponentials | steps: cases of cos x; value of f; continuity}]
Consider the function {\scriptsize$f(x)=\displaystyle\lim_{n\to\infty}\dfrac{(1+\cos x)^n+5\log_ex}{2+(1+\cos x)^n}$}, then
\begin{qoptions}
\opt{A}{$f(x)$ is continuous at positive odd integral multiple of $\pi$}
\opt{B}{$f(x)$ is discontinuous at positive even integral multiple of $\pi$}
\opt{C}{$f(x)$ is discontinuous at positive odd integral multiple of $\dfrac\pi2$}
\opt{D}{$f(x)$ is continuous at positive even integral multiples of $\dfrac\pi2$}
\end{qoptions}
\end{bankq}
\begin{bankq}[id=TNAS02-P1-MATH-S2-Q07, source={TNAS - Test 02 (07 May 2023)}, paper=1, qno=7, type=multi, subtopic={Standard limits}, answer={A,C}, answer_src=key, difficulty=6, diff_parts={K1 S2 I2 C1 F1}, diff_note={concepts: limit of a sum of logarithms | steps: combine; expansion of tangent; sum; relation of the two terms}]
If {\footnotesize$P_n=\displaystyle\lim_{x\to0^+}$}\par\noindent\hfill{\scriptsize$\left(\log(\tan x)+\log(\tan2x)+\log(\tan3x)+\ldots+\log(\tan nx)-\log x^n\right)$}\hfill\null\par\par\noindent and $P_{256}-P_{255}=kP_m$, (where $k,m,n\in N$), then-
\begin{qoptions}
\opt{A}{$k=8$}
\opt{B}{$k=6$}
\opt{C}{$m=2$}
\opt{D}{$m=4$}
\end{qoptions}
\end{bankq}
\begin{bankq}[id=TNAS02-P1-MATH-S4-Q09, source={TNAS - Test 02 (07 May 2023)}, paper=1, qno=9, type=numerical, subtopic={Standard limits}, answer={16}, answer_src=key, difficulty=6, diff_parts={K1 S2 I2 C1 F1}, diff_note={concepts: limit at minus infinity | steps: rationalise; leading terms; value}]
If $\displaystyle\lim_{x\to-\infty}\left(\sqrt{x^2-\dfrac{4x}{3}}+x+2\right)$ is equal to $K$ then value of $6K$ is-
\end{bankq}
\begin{bankq}[id=TNAS02-P2-MATH-S4-Q08, source={TNAS - Test 02 (07 May 2023)}, paper=2, qno=8, type=numerical, subtopic={Standard limits}, answer={1}, answer_src=key, difficulty=6, diff_parts={K1 S2 I2 C1 F1}, diff_note={concepts: right hand derivative | steps: expand for small fractional part; limit; absolute value}]
If $f(x)=\begin{cases}\dfrac{\sin^{-1}(\{x\}^2+\{x\}^3)-x}{\{x\}}, & x\neq0\\ -1, & x=0\end{cases}$, then the absolute value of RHD of $f(x)$ at $x=0$, is\\
(where $\{.\}$ denotes fractional part function)
\end{bankq}
\begin{bankq}[id=TNAS05-P1-MATH-S1-Q01, source={TNAS - Test 05 (09 Jul 2023)}, paper=1, qno=1, type=single, subtopic={L'Hopital}, answer={A}, answer_src=key, difficulty=3, diff_parts={K1 S2 I1 C0 F0}, diff_note={concepts: limit of an integral of a function | steps: substitute the integral; L'Hopital; evaluate}]
Let $f:[0,1]\to R$ be a derivable function with $f(0)=0$, then $\displaystyle\lim_{n\to\infty}n^2\int_0^{1/n}f(t)\,dt$ is equal to-
\begin{qoptions}
\opt{A}{$\dfrac{f'(0)}{2}$}
\opt{B}{$f'(0)$}
\opt{C}{$2f'(0)$}
\opt{D}{0}
\end{qoptions}
\end{bankq}
\begin{bankq}[id=TNAS07-P1-MATH-S1-Q04, source={TNAS - Test 07 (20 Aug 2023)}, paper=1, qno=4, type=single, subtopic={Standard limits}, answer={B}, answer_src=key, difficulty=7, diff_parts={K2 S3 I2 C1 F0}, diff_note={concepts: limit with fractional part, polynomial roots | steps: behaviour of the numerator; factors of the denominator; conditions at one and two; polynomial; sum of coefficients}]
Consider $f(x)=\dfrac{1-\cos\{x\}}{\left(x^4+ax^3+bx^2+cx+d\right)^2}$ (where $\{.\}$ fractional part of function). If $\displaystyle\lim_{x\to1^+}f(x)$ and $\displaystyle\lim_{x\to2^+}f(x)$ both exist and are non zero then $a+b+c+d$ is equal to
\begin{qoptions}
\opt{A}{2}
\opt{B}{$-1$}
\opt{C}{0}
\opt{D}{5}
\end{qoptions}
\end{bankq}
\begin{bankq}[id=TNAS07-P1-MATH-S1-Q08, source={TNAS - Test 07 (20 Aug 2023)}, paper=1, qno=8, type=single, subtopic={Standard limits}, answer={C}, answer_src=key, difficulty=7, diff_parts={K1 S3 I2 C2 F0}, diff_note={concepts: limit of a logarithm to a trigonometric base | steps: change of base; expansions; limit of the ratio; square; cube}]
$\displaystyle\lim_{x\to0}\left(2+\log^2_{\sec\left(\frac x2\right)}\cos\left(\dfrac x3\right)\right)^3$ is equal to
\begin{qoptions}
\opt{A}{$\left(\dfrac{146}{81}\right)^3$}
\opt{B}{$\left(\dfrac{70}{27}\right)^3$}
\opt{C}{$\left(\dfrac{178}{81}\right)^3$}
\opt{D}{$\left(\dfrac{34}9\right)^3$}
\end{qoptions}
\end{bankq}
\begin{bankq}[id=TNAS07-P1-MATH-S2-Q07, source={TNAS - Test 07 (20 Aug 2023)}, paper=1, qno=7, type=numerical, subtopic={Expansions}, answer={3}, answer_src=solved, working={With $t=x^{-1/2}$: $\ln\sin t=\ln t-\dfrac{t^2}6$, so the bracket is $\dfrac1{3x}$ and $L=\dfrac13$, $\dfrac1L=3$}, difficulty=8, diff_parts={K1 S3 I2 C2 F1}, diff_note={concepts: limit at infinity with a logarithm and a sine | steps: substitute reciprocal; expand the sine; logarithm series; limit; reciprocal}]
If $L=\displaystyle\lim_{x\to\infty}x\ell n\left(\dfrac1x\right)-2x\ell n\left(\sin\dfrac1{\sqrt x}\right)$ then $\dfrac1L$ is equal to
\end{bankq}
\begin{bankq}[id=TNAS07-P2-MATH-S2-Q07, source={TNAS - Test 07 (20 Aug 2023)}, paper=2, qno=7, type=single, subtopic={Standard limits}, answer={A}, answer_src=key, difficulty=7, diff_parts={K1 S3 I2 C2 F0}, diff_note={concepts: four limit and function items matched | steps: change of base; limit at infinity; range condition; implicit differentiation; match}]
{\small\setlength{\tabcolsep}{2pt}\begin{tabular}{|c|p{0.58\linewidth}|c|c|}
\hline
\multicolumn{2}{|c|}{\textbf{List-I}} & \multicolumn{2}{c|}{\textbf{List-II}}\\ \hline
(I) & Value of $\lim\limits_{x\to0}\log_{\frac{x^2+1}{x+1}}\left(\sqrt{x^2+1}-x\right)$ is less than or equal to & (P) & $-1$\\ \hline
(II) & The value of $\lim\limits_{x\to\infty}\dfrac{2\left(3^{\sin x}+2x+1\right)}{\sin x+\sqrt{x^2+1}}$ is greater than or equal to & (Q) & $4$\\ \hline
(III) & Let $f:R\to\left[0,\dfrac\pi2\right)$ be defined by $f(x)=\tan^{-1}\left(3x^2+6x+a\right)$. If $f(x)$ is an onto function then the value of $a$ can be & (R) & $3$\\ \hline
(IV) & Let $x^5+10x^3y^2+15x^2y^3+3x^4y+y^5=0$ and $\dfrac{dy}{dx}=\dfrac{\lambda y}x$, then $\lambda$ is equal to & (S) & $5$\\ \hline
 & & (T) & $1$\\ \hline
\end{tabular}}\par
\begin{qoptions}
\opt{A}{$\mathrm I\to Q,R,S,T;\allowbreak \mathrm{II}\to P,Q,R,T;\allowbreak \mathrm{III}\to R;\allowbreak \mathrm{IV}\to T$}
\opt{B}{$\mathrm I\to P,Q,R,T;\allowbreak \mathrm{II}\to Q,R,S,T;\allowbreak \mathrm{III}\to R;\allowbreak \mathrm{IV}\to T$}
\opt{C}{$\mathrm I\to Q,R,S,T;\allowbreak \mathrm{II}\to R;\allowbreak \mathrm{III}\to P,Q,R,T;\allowbreak \mathrm{IV}\to T$}
\opt{D}{$\mathrm I\to Q,R,S,T;\allowbreak \mathrm{II}\to P,Q,R,T;\allowbreak \mathrm{III}\to T;\allowbreak \mathrm{IV}\to R$}
\end{qoptions}
\end{bankq}
\begin{bankq}[id=TNAS07-P2-MATH-S3-Q02, source={TNAS - Test 07 (20 Aug 2023)}, paper=2, qno=2, type=numerical, subtopic={Standard limits}, answer={1.00}, answer_src=key, difficulty=7, diff_parts={K1 S2 I2 C2 F1}, diff_note={concepts: limit of two indeterminate forms | steps: logarithm of each term; limit of each; add}]
$\displaystyle\lim_{x\to0^+}(x\cos x)^x+(x\sin x)^{1/x}$ is
\end{bankq}

\section*{Internal Test papers}
\begin{bankq}[id=INT01-P1-MATH-S1-Q01, source={Internal Test - Test 01 (05 Aug 2020)}, paper=1, qno=1, type=single, subtopic={Standard limits}, answer={C}, answer_src=key, difficulty=6, diff_parts={K1 S2 I2 C2 F0}, diff_note={concepts: functional equation of the tangent addition type | steps: identify the function as an inverse tangent; use the limit; constant; evaluate}]
If $f(x)+f(y)=f\left(\dfrac{x+y}{1-xy}\right)$ for all $x,y\in\mathbb R$ and $xy\ne1$, and $\displaystyle\lim_{x\to0}\dfrac{f(x)}{x}=2$, the value of $f\left(\dfrac1{\sqrt3}\right)$ is
\begin{qoptions}
\opt{A}{$1$}
\opt{B}{$\dfrac13$}
\opt{C}{$\dfrac\pi3$}
\opt{D}{$\dfrac\pi2$}
\end{qoptions}
\end{bankq}
\begin{bankq}[id=INT01-P1-MATH-S1-Q03, source={Internal Test - Test 01 (05 Aug 2020)}, paper=1, qno=3, type=single, subtopic={Expansions}, answer={D}, answer_src=key, difficulty=6, diff_parts={K1 S2 I2 C2 F0}, diff_note={concepts: limit using the function with a repeated factor | steps: expansion of g near zero; third derivative of f; limit; lambda}]
Let $f(x)$ and $g(x)$ be two functions such that $f(x)=x^2g(x)$, $\displaystyle\lim_{x\to0}g(x)\csc x=\dfrac{f'''(0)}{\lambda}$.\par
The value of $\lambda$ is
\begin{qoptions}
\opt{A}{$2$}
\opt{B}{$3$}
\opt{C}{$4$}
\opt{D}{$6$}
\end{qoptions}
\end{bankq}
\begin{bankq}[id=INT02-P1-MATH-S1-Q01, source={Internal Test - Test 02 (07 Jul 2021)}, paper=1, qno=1, type=single, subtopic={Standard limits}, answer={D}, answer_src=key, difficulty=5, diff_parts={K1 S2 I2 C1 F0}, diff_note={concepts: limit of a ratio of a decreasing function | steps: monotonic bounds; squeeze the ratio; limit}]
Let $f:(1,\infty)\to(0,\infty)$ be continuous strictly decreasing function with $\displaystyle\lim_{x\to\infty}\dfrac{f(4x)}{f(8x)}=1$. Then $\displaystyle\lim_{x\to\infty}\dfrac{f(6x)}{f(8x)}=$
\begin{qoptions}
\opt{A}{$\dfrac12$}
\opt{B}{$\dfrac23$}
\opt{C}{$\dfrac34$}
\opt{D}{$1$}
\end{qoptions}
\end{bankq}
\begin{bankq}[id=INT02-P1-MATH-S2-Q07, source={Internal Test - Test 02 (07 Jul 2021)}, paper=1, qno=7, type=multi, subtopic={Standard limits}, answer={B,D}, answer_src=key, difficulty=7, diff_parts={K1 S2 I2 C2 F1}, diff_note={concepts: limit with fractional part and conjugate powers | steps: conjugate sum is an integer; fractional part; cosine of the integer part; parity cases}]
$\displaystyle\lim_{n\to\infty}\left\{\left(\sqrt5+2\right)^n\right\}\cdot\cos\left(\pi\left[\left(\sqrt5+2\right)^n\right]\right)$ is equal to\par
(where $n\in I^+$ and $[\bullet],\{\bullet\}$ denotes greatest integer function and fractional part function respectively)
\begin{qoptions}
\opt{A}{$0$ when $n$ is even}
\opt{B}{$-1$ when $n$ is even}
\opt{C}{$-1$ when $n$ is odd}
\opt{D}{$0$ when $n$ is odd}
\end{qoptions}
\end{bankq}
\begin{bankq}[id=INT02-P1-MATH-S3-Q14, source={Internal Test - Test 02 (07 Jul 2021)}, paper=1, qno=14, type=numerical, subtopic={Standard limits}, answer={10}, answer_src=key, difficulty=7, diff_parts={K1 S2 I2 C2 F1}, diff_note={concepts: composition of limits with continuous functions | steps: value of f at three; value of g at the limit; limit of tangent}]
Let $f(x)\ \&\ g(x)$ be continuous function for all $x\in R$. If $\displaystyle\lim_{x\to3}\dfrac{f(x)}{x-3}=\lim_{x\to f(3)}\dfrac{g(x)}{x^2}=5$ then $\displaystyle\lim_{x\to g(f(3))}\dfrac{\tan10x}{x}=$
\end{bankq}
\begin{bankq}[id=INT02-P1-MATH-S3-Q18, source={Internal Test - Test 02 (07 Jul 2021)}, paper=1, qno=18, type=numerical, subtopic={Standard limits}, answer={3}, answer_src=key, difficulty=7, diff_parts={K1 S2 I2 C2 F1}, diff_note={concepts: limit using a derivative of a polynomial | steps: derivative at one; expansion of each factor; limit; scale}]
Let $f(x)=3x^{10}-7x^8+5x^6-21x^3+3x^2-7$\par
$265\left(\displaystyle\lim_{h\to0}\dfrac{h^4+3h^2}{\left(f(1-h)-f(1)\right)\sin5h}\right)=$
\end{bankq}
\end{document}
